
\providecommand{\bk}{\mathbb{k}}
\chapter{Nilpotent completion for filtered chiral bar--cobar}
\label{app:nilpotent-completion}


A filtered chiral algebra can have finite bar words and still fail to
define a discrete bar--cobar comparison: a fixed output window may
receive infinitely many OPE-mode contributions, curvature terms may
survive in filtration degree~$0$, and the universal twisting
cochain must be evaluated in a completed convolution algebra.  The
completion used here is not a substitute for Koszul duality.  It is the
topology in which the bar coalgebra, the cobar differential, and the
Maurer--Cartan equation are defined by finite-window limits.

\begin{remark}[Three bar surfaces]
\label{rem:nilpotent-completion-dependency}
There are three distinct objects.
\[
\Omega_X\bar B_X(A)\longrightarrow A
\]
is the bar--cobar counit; when it is a quasi-isomorphism it is
inversion.  The completed object
\[
\widehat{\bar B}_X(A)=\varprojlim_N \bar B_X(A_{\le N})
\]
is a pro-conilpotent bar coalgebra.  The Verdier--Koszul algebra is the
separate construction
\[
A^!_\infty:=\mathbb D_{\Ran}\widehat{\bar B}_X(A),
\]
under the dualizability and continuity hypotheses for Verdier duality
on the Ran space.  The cobar counit belongs to the reconstruction lane:
$\Omega_X\widehat{\bar B}_X(A)\simeq A$ recovers~$A$, while
$\mathbb D_{\Ran}\widehat{\bar B}_X(A)$ defines the Verdier--Koszul
algebra.
\end{remark}

\section{Finite windows and completion}

\subsection{Filtered non-quadratic algebras}

\begin{observation}[Filtered bar words]\label{obs:quadratic-limitation}
Classical quadratic Koszul duality sees generators and quadratic
relations.  Chiral algebras with central terms, normally ordered
composites, or infinitely many positive-weight generators carry more
data than this finite quadratic surface.  The Virasoro OPE
\[
T(z)T(w)\sim
\frac{c/2}{(z-w)^4}+\frac{2T(w)}{(z-w)^2}
 +\frac{\partial T(w)}{z-w}
\]
is already curved: the quartic pole contributes a scalar curvature
term in the bar differential.  Principal $\mathcal W_N$ algebras add
composite fields.  Positive towers such as $\mathcal W_{1+\infty}$ or
RTT/Yangian towers require a weight-completed inverse limit of finite
windows.

The deconcatenation coalgebra remains conilpotent on every finite bar
word.  Completion is needed for a different reason: the bar
differential, the cobar differential, and the universal twisting
morphism must be continuous for the chosen filtration, and the
inverse-limit cohomology must satisfy Mittag--Leffler.
\end{observation}

\subsection{The finite-window topology}

\begin{definition}[Finite-window completion]\label{idea:I-adic}
Let $A$ be an augmented filtered chiral $A_\infty$-algebra on a curve
$X$ with
\[
A=F^0A\supset F^1A=\bar A\supset F^2A\supset\cdots,
\qquad
A\simeq\varprojlim_N A_{\le N},
\qquad
A_{\le N}:=A/F^{N+1}A .
\]
The filtration is a \emph{strong finite-window filtration} if
\[
\mu_r(F^{i_1}A,\ldots,F^{i_r}A)
\subset F^{i_1+\cdots+i_r}A
\tag{\ast}
\label{eq:nc-strong-filtration}
\]
for every chiral $A_\infty$ operation $\mu_r$, and every finite window
$A_{\le N}$ lies in the finite-type bar--cobar regime.
The completed bar coalgebra is
\[
\widehat{\bar B}_X(A)
:=\varprojlim_N \bar B_X(A_{\le N}).
\]
For an honest augmentation-ideal filtration this is the usual
$I$-adic completion; for a conformal-weight tower it is the
weight-window completion.
\end{definition}

\section{The conilpotent filtration}

\subsection{Definition and properties}

\begin{definition}[Conilpotent filtration]\label{def:conilpotent-filtration}
For a coaugmented coalgebra $C$ with counit $\epsilon: C \to k$, coaugmentation $\eta: k \to C$, and coproduct $\Delta: C \to C \otimes C$, the \emph{coradical filtration} (or conilpotent filtration) is:
\begin{align*}
F_0C &= \operatorname{im}(\eta) = k \quad (\text{coaugmentation}) \\
F_nC &= \{c \in C : \Delta(c) \in F_{n-1}C \otimes C + C \otimes F_{n-1}C\} \quad (n \geq 1)
\end{align*}
Equivalently, using the reduced coproduct $\bar{\Delta}: \ker(\epsilon) \to \ker(\epsilon) \otimes \ker(\epsilon)$ and its iterates $\bar{\Delta}^{(n)}: \ker(\epsilon) \to \ker(\epsilon)^{\otimes(n+1)}$:
\[
F_nC = k \oplus \ker(\bar{\Delta}^{(n)}).
\]
In particular, $F_1C = k \oplus \ker(\bar{\Delta})$ consists of the scalars and the primitive elements.

This is an \emph{ascending} filtration: $F_0C \subset F_1C \subset F_2C \subset \cdots$.
A coalgebra is conilpotent if $\bigcup_n F_nC = C$, equivalently, if every element of $\ker(\epsilon)$ is annihilated by a sufficiently high iterate of the reduced coproduct: $\ker(\epsilon) = \bigcup_n \ker(\bar{\Delta}^{(n)})$.
\end{definition}

\begin{proposition}[Collision trees and coradical degree; \ClaimStatusProvedHere]\label{prop:geom-conilpotent}
For the geometric bar model
\[
\bar B^{\mathrm{geom}}_r(\mathcal A)
=
\Gamma(\overline C_r(X),
\mathcal A^{\boxtimes r}\otimes\Omega^*_{\log}),
\]
the coradical filtration is the finite bar-length filtration:
\[
F_k\bar B^{\mathrm{geom}}(\mathcal A)
=
\bk\oplus\bigoplus_{1\le r\le k}
\bar B^{\mathrm{geom}}_r(\mathcal A).
\]
The Fulton--MacPherson collision-tree filtration is a geometric
refinement of this filtration, not a direct-sum decomposition by
support on boundary strata.  If $G_k$ denotes the subspace of log
forms whose iterated residue expansion has no tree of depth $>k$,
then
\[
\bar\Delta(G_k)\subset
\sum_{i+j\le k-1}G_i\otimes G_j.
\]
Thus collision depth is a geometric witness for conilpotence, while
the coradical degree itself is measured by iterated deconcatenation.
\end{proposition}

\begin{proof}
On the tensor coalgebra $T^c(s^{-1}\bar A)$ the reduced coproduct is
deconcatenation:
\[
\bar\Delta[a_1|\cdots|a_r]
=
\sum_{i=1}^{r-1}
[a_1|\cdots|a_i]\otimes[a_{i+1}|\cdots|a_r].
\]
After $r$ reduced coproducts every word of length~$r$ has vanished.
Hence the coradical filtration is exactly the filtration by word
length, with the coaugmentation in degree~$0$.

The compactification $\overline C_r(X)$ is stratified by rooted
collision trees.  An iterated logarithmic residue chooses a chain of
boundary divisors, equivalently a nested collision tree.  Cutting such
a tree along an internal edge is the geometric incarnation of a
deconcatenation cut.  Each cut shortens the maximal root-to-leaf chain
in the remaining factors, so the displayed inclusion for $G_k$ follows.
Every collision tree is finite, and every bar word has finite length;
therefore the filtration is ascending and exhaustive.  No assertion is
made that logarithmic forms split as a direct sum of forms supported on
strata.
\end{proof}

\subsection{Convergence by finite windows}

\begin{theorem}[Finite-window convergence; \ClaimStatusProvedHere]\label{thm:completion-convergence}
Let $A$ be a separated complete augmented filtered chiral
$A_\infty$-algebra with a strong finite-window filtration
\textup{(}Definition~\textup{\ref{idea:I-adic}}\textup{)}, and assume
that the finite-stage maps induced on bar complexes are strict
epimorphisms of complexes on every finite bar-length window.  Then:
\begin{enumerate}
\item the inverse limit
 $\widehat{\bar B}_X(A)=\varprojlim_N\bar B_X(A_{\le N})$
 is a separated complete pro-conilpotent curved dg chiral coalgebra;
\item the bar differential, coproduct, curvature, and universal
 twisting morphism are continuous for the finite-window topology;
\item the completed convolution algebra
 \[
 \widehat{\operatorname{Conv}}(\widehat{\bar B}_X(A),A)
 :=
 \varprojlim_N
 \operatorname{Conv}(\bar B_X(A_{\le N}),A_{\le N})
 \]
 contains the compatible limit
 $\widehat\tau=\varprojlim_N\tau_N$, and
 $\widehat\tau$ satisfies the Maurer--Cartan equation
 \[
 \partial\widehat\tau+\widehat\tau\star\widehat\tau=0;
 \]
\item for each cohomological degree~$m$ there is a Milnor exact
 sequence
 \[
 0\to
 \varprojlim\nolimits_N^1 H^{m-1}(\bar B_X(A_{\le N}))
 \to
 H^m(\widehat{\bar B}_X(A))
 \to
 \varprojlim_N H^m(\bar B_X(A_{\le N}))
 \to0;
 \]
\item if, in degree~$m$, the finite-stage cohomology system
 $\{H^m(\bar B_X(A_{\le N}))\}_N$ is strict Mittag--Leffler, then
 \[
 H^m(\widehat{\bar B}_X(A))
 \cong
 \varprojlim_N H^m(\bar B_X(A_{\le N})).
 \]
\end{enumerate}
The strong finite-window multiplicativity condition of
Definition~\ref{idea:I-adic} implies the needed finiteness of the
Maurer--Cartan equation in every window: modulo $F^{N+1}A$, only
operations $\mu_r$ with $r\le N$ can contribute on reduced inputs.
\end{theorem}

\begin{proof}
The quotient maps $A_{\le N+1}\twoheadrightarrow A_{\le N}$ are
strict morphisms of curved chiral $A_\infty$-algebras.  Functoriality
of the finite-stage bar construction gives morphisms of curved dg
chiral coalgebras
\[
\bar B_X(A_{\le N+1})\longrightarrow\bar B_X(A_{\le N})
\]
compatible with coproduct, coaugmentation, differential, and curvature.
The strict epimorphism hypothesis makes this a strict inverse system
of complexes, so the inverse limit carries these structures
componentwise in the category of separated complete complexes.
Each finite quotient is conilpotent by bar length; the limit is
pro-conilpotent and separated because $A\simeq\varprojlim_N A_{\le N}$.

For reduced inputs, the strong finite-window condition gives
$\mu_r(\bar A^{\otimes r})\subset F^rA$.  Hence modulo $F^{N+1}A$ all
terms with $r>N$ vanish.  The universal finite-stage twisting
morphisms $\tau_N$ satisfy
$\partial\tau_N+\tau_N\star\tau_N=0$ and are compatible with the
transition maps; their inverse limit is the displayed completed
Maurer--Cartan element.  Continuity is precisely the statement that
every component is detected after projection to some finite window.

The cohomology statement is Milnor's exact sequence for a strict tower
of complexes:
\[
0\to
\varprojlim\nolimits_N^1 H^{m-1}(\bar B_X(A_{\le N}))
\to
H^m(\widehat{\bar B}_X(A))
\to
\varprojlim_N H^m(\bar B_X(A_{\le N}))
\to0 .
\]
Strict Mittag--Leffler kills the derived limit term.  This is the only
place where exactness of inverse limits is used.
\end{proof}

\begin{remark}[Polynomial growth is not a convergence theorem]\label{rem:completion-convergence-window}
Finite generation, polynomial bounds on OPE coefficients, and
finite-dimensional Hochschild cohomology do not by themselves imply
eventual vanishing of
$H^2(A,F^N/F^{N+1})$.  The convergence statement used in this chapter
is the finite-window theorem above: either the strong filtration axiom
gives a degree cutoff, or a separate weight-cutoff/Mittag--Leffler
criterion verifies the inverse-limit comparison.
\end{remark}

\section{Completion and the Verdier--Koszul lane}

\subsection{The main construction}

\begin{construction}[Completed bar coalgebra and Koszul dual lane]\label{const:completed-koszul}
For a complete filtered algebra~$A$, set
\[
\widehat B_X(A):=\widehat{\bar B}_X(A)
=\varprojlim_N\bar B_X(A_{\le N}).
\]
This is a completed dg coalgebra.  The coalgebraic Koszul object is
the bar cohomology
\[
A^{\mathrm i}_{\mathrm{cpl}}:=H^\bullet(\widehat B_X(A))
\]
when that cohomology is a usable completed coalgebra.  The
Verdier--Koszul algebra is
\[
A^!_\infty:=\mathbb D_{\Ran}\widehat B_X(A),
\]
under the topological finiteness hypotheses that make
$\mathbb D_{\Ran}$ continuous on the completed bar object.

The completed cobar used for inversion is a different object:
\[
\widehat\Omega_X(\widehat B_X(A))
:=\varprojlim_N\Omega_X(\bar B_X(A_{\le N})).
\]
\end{construction}

\begin{theorem}[Completed finite-window bar--cobar inversion; \ClaimStatusProvedHere]\label{thm:completed-bar-cobar}
Let $A$ be as in Theorem~\textup{\ref{thm:completion-convergence}},
and assume the finite-stage counits
\[
\epsilon_N\colon \Omega_X\bar B_X(A_{\le N})\longrightarrow A_{\le N}
\]
are compatible quasi-isomorphisms in the square-zero total model.
Assume also that the cone tower
$K_N=\operatorname{Cone}(\epsilon_N)$ is strict Mittag--Leffler
\textup{(}for example, the transition maps
$K_{N+1}\twoheadrightarrow K_N$ are degreewise surjective\textup{)}.
Then the completed counit
\[
\widehat\epsilon\colon
\widehat\Omega_X(\widehat{\bar B}_X(A))
=\varprojlim_N\Omega_X\bar B_X(A_{\le N})
\longrightarrow
\varprojlim_N A_{\le N}=A
\]
is a quasi-isomorphism.  In a genuinely curved CDG realization the
same finite-window construction gives a coacyclic cone in the relevant
second-kind coderived category, not an ordinary raw-chain
quasi-isomorphism.
\end{theorem}

\begin{proof}
The maps $\epsilon_N$ form a morphism of inverse systems.  Let
$K_N=\operatorname{Cone}(\epsilon_N)$.  Each $K_N$ is acyclic in the
square-zero total model by the finite-stage bar--cobar theorem.  The
strict Mittag--Leffler hypothesis on the cone tower is part of the
completed statement; it is not a formal consequence of acyclicity of
the individual cones.  The Milnor
sequence gives
\[
0\to\varprojlim\nolimits_N^1H^{m-1}(K_N)
\to H^m(\varprojlim_NK_N)
\to \varprojlim_NH^m(K_N)\to0.
\]
The right term vanishes because every $K_N$ is acyclic; the left term
vanishes by Mittag--Leffler.  Hence
$\operatorname{Cone}(\widehat\epsilon)$ is acyclic, which is the
displayed quasi-isomorphism.  In the curved CDG ambient the same
inverse-limit cone is interpreted in Positselski's second-kind
category; acyclicity is replaced by coacyclicity.
\end{proof}

\begin{remark}[Finite-stage compatibility]\label{rem:completed-bar-cobar-window}
The theorem uses finite-stage inversion as an input.  The statement
does not assert that an arbitrary quotient of an arbitrary filtered
algebra lies on the Koszul locus; it asserts that a tower whose
quotients already carry the finite-stage bar--cobar quasi-isomorphisms
inherits the completed quasi-isomorphism by Milnor exactness.
\end{remark}

\subsection{Essential image}

\begin{theorem}[Essential image of finite-window bar towers; \ClaimStatusProvedHere]\label{thm:koszul-dual-characterization}
A separated complete curved dg chiral coalgebra
$\widehat C$ is isomorphic to $\widehat{\bar B}_X(A)$ for a strong
finite-window algebra~$A$ if and only if it admits a presentation
\[
\widehat C\simeq\varprojlim_N C_N
\]
with the following properties:
\begin{enumerate}
\item each $C_N$ is a finite-type conilpotent curved dg chiral
 coalgebra in the finite-stage essential image, say
 $C_N\simeq\bar B_X(A_{\le N})$;
\item the transition maps $C_{N+1}\to C_N$ are induced by strict
 quotient maps $A_{\le N+1}\twoheadrightarrow A_{\le N}$;
\item the algebras $A_{\le N}$ assemble to a separated complete
algebra $A=\varprojlim_NA_{\le N}$ whose operations satisfy the
 strong finite-window multiplicativity condition of
 Definition~\ref{idea:I-adic}; and
\item the finite-stage twisting morphisms and counits are compatible
 with the tower.
\end{enumerate}
\end{theorem}

\begin{proof}
If $\widehat C=\widehat{\bar B}_X(A)$, take
$C_N=\bar B_X(A_{\le N})$.  The four conditions are exactly the
finite-window construction of Definition~\ref{idea:I-adic} and
Theorem~\ref{thm:completion-convergence}.

Conversely, assume such a presentation.  The strict quotient maps
$A_{\le N+1}\twoheadrightarrow A_{\le N}$ assemble to the separated
complete algebra $A=\varprojlim_NA_{\le N}$.  Compatibility of the
finite-stage bar functor identifies the inverse system
$\{C_N\}_N$ with $\{\bar B_X(A_{\le N})\}_N$.  Taking inverse limits
gives
\[
\widehat C
\simeq
\varprojlim_N C_N
\simeq
\varprojlim_N\bar B_X(A_{\le N})
=
\widehat{\bar B}_X(A).
\]
\end{proof}

\begin{remark}[Spectral sequences do not define the essential image]\label{rem:koszul-dual-characterization-window}
Degeneration of an associated-graded spectral sequence is useful for
proving that a given finite-window tower lies on the Koszul locus, but
it is not an intrinsic characterization of completed bar coalgebras.
The essential-image datum is the compatible finite-stage bar tower
itself, together with the continuous twisting morphisms.
\end{remark}

\section{Connection to Beilinson--Drinfeld}

\subsection{Filtered chiral algebras}

\begin{remark}[BD framework]\label{rem:BD-filtered}
Beilinson--Drinfeld develop filtered chiral algebras with
augmentation-ideal topology.  The basic operations are:

\begin{itemize}
\item filtration by the augmentation ideal~$I$;
\item completion
 $\widehat{\mathcal{A}}=\varprojlim_n\mathcal{A}/I^n$;
\item associated graded
 $\operatorname{gr}(\mathcal{A})=\bigoplus_n I^n/I^{n+1}$.
\end{itemize}
When the finite-window filtration is the augmentation-ideal
filtration, Definition~\ref{idea:I-adic} is exactly this completion
applied to the bar coalgebra:
\[
\widehat{\bar B}_X(A)=\varprojlim_n\bar B_X(A/I^{n+1}).
\]
\end{remark}

\subsection{Chiral homology completion}

\begin{theorem}[BD chiral homology \cite{BD04}; \ClaimStatusProvedElsewhere]\label{thm:BD-chiral-homology}
Let $A$ be an $I$-adically complete filtered chiral algebra on a
smooth curve~$X$.  Assume the inverse system
$\{H_*^{\mathrm{ch}}(X,A/I^{n+1})\}_n$ is Mittag--Leffler.  Then
completed chiral homology is computed by the completed bar coalgebra:
\[
\widehat H_*^{\mathrm{ch}}(X,A)
:=
\varprojlim_n H_*^{\mathrm{ch}}(X,A/I^{n+1})
\cong
H_*\!\left(\varprojlim_n\bar B_X(A/I^{n+1})\right).
\]
\end{theorem}

\section{Perturbative filtrations}

\subsection{Perturbative expansion}

\begin{interpretation}[Completion as perturbative expansion]\label{interp:completion-renorm}
In perturbative models whose Feynman graphs are identified with the
bar construction, a finite window $\bar B(A_{\le N})$ records the graph
weights visible below the chosen filtration bound.  The completed
object
\[
\widehat{\bar B}_X(A)=\varprojlim_N\bar B_X(A_{\le N})
\]
is the formal all-window object.  Theorem~\ref{thm:completion-convergence}
is the algebraic convergence statement: finite graph windows assemble
only when the differential and twisting cochain are continuous and the
cohomology tower is Mittag--Leffler.
\end{interpretation}

\subsection{Feynman diagram expansion}

\begin{remark}[Feynman diagrams and conilpotent degree]\label{rem:feynman-conilpotent}
The conilpotent filtration levels $F^n/F^{n+1}$ correspond, under the Feynman diagram identification of Chapter~\ref{ch:v1-feynman}, to graphs with exactly~$n$ internal edges. In particular: $F^0/F^1$ consists of vacuum contributions (no internal edges), $F^1/F^2$ of tree-level propagators (single internal edge), and $F^n/F^{n+1}$ of diagrams with~$n$ internal edges. The symmetry factors arising in the perturbative expansion are the orders of the automorphism groups of the corresponding graphs. The completion $\widehat{\bar{B}}(\mathcal{A}) = \varprojlim_n \bar{B}(\mathcal{A})/I^n$ amounts to summing over all graph topologies, a formal power series that requires the $I$-adic topology to converge.
\end{remark}

\section{Finite-window examples}

\subsection{Example 1: Virasoro algebra}

\begin{example}[Virasoro finite-window completion]\label{ex:virasoro-completion}
The Virasoro OPE is
\[
T(z)T(w)\sim
\frac{c/2}{(z-w)^4}
+\frac{2T(w)}{(z-w)^2}
+\frac{\partial T(w)}{z-w}.
\]
The quartic pole contributes the scalar curvature direction
$m_0\propto(c/2)\mathbf 1$ in the curved bar model.  On the
positive-energy state space write
\[
(\mathrm{Vir}_c)_{\le N}:=\bigoplus_{0\le h\le N}
(\mathrm{Vir}_c)_h .
\]
Each finite window is finite-dimensional by the VOA finite-energy
axiom, and the completed bar coalgebra is
\[
\widehat{\bar B}(\mathrm{Vir}_c)
=
\varprojlim_N
\bar B\bigl((\mathrm{Vir}_c)_{\le N}\bigr).
\]
This object is the completed bar coalgebra.  The same-family shadow
$\mathrm{Vir}_{26-c}$ belongs to the Verdier--Koszul lane; it is not
the bar coalgebra itself and is not produced by the counit
$\Omega\widehat{\bar B}(\mathrm{Vir}_c)\to\mathrm{Vir}_c$.
\end{example}

\subsection{\texorpdfstring{Example 2: $W_3$ algebra}{Example 2: W-3 algebra}}

\begin{example}[\texorpdfstring{$\mathcal W_3$}{W3} finite windows]\label{ex:w3-completion}
The algebra $\mathcal W_3$ has strong generators $T$ and~$W$ of spins
$2$ and~$3$, and the $W$--$W$ OPE contains a sixth-order scalar pole.
The completion is the conformal-weight tower
\[
\widehat{\bar B}(\mathcal W_3)
=
\varprojlim_N
\bar B\bigl((\mathcal W_3)_{\le N}\bigr).
\]
The convergence input is not the bare assertion that the structure
constants have polynomial growth.  The finite-window proof requires
the operations to preserve the chosen filtration, finite-stage
bar--cobar inversion on each $(\mathcal W_3)_{\le N}$, and
Mittag--Leffler stabilization of the resulting bar cohomology tower.
\end{example}

\subsection{Example 3: Heisenberg (quadratic case for comparison)}

\begin{example}[Heisenberg: discrete conilpotent model]\label{ex:heisenberg-no-completion}
The Heisenberg algebra is quadratic:
\[a(z)a(w) \sim \frac{k}{(z-w)^2}\]

At genus~$0$ the reduced bar coalgebra is conilpotent by bar length:
for a word of length~$n$, the $(n+1)$-fold reduced coproduct vanishes.
The quadratic Koszul model is finite in each finite input window, and
the bar--cobar counit is defined in the discrete conilpotent category.
Completing the conformal-weight tower changes no finite bar word; it
only records the same model as a separated inverse limit.
\end{example}

\section{Computational methods}

\subsection{Computing the completion}

\begin{remark}[Computing a finite window]\label{rem:compute-completion}
Given a filtered chiral algebra~$A$ and a window~$N$, the quotient
$A_{\le N}=A/F^{N+1}A$ determines the finite-stage bar coalgebra
$\bar B_X(A_{\le N})$.
\begin{enumerate}
\item \emph{Basis.} Keep fields whose filtration degree is at most~$N$.
\item \emph{Differential.} Compute collision residues from the OPE and
 discard every output in $F^{N+1}A$.
\item \emph{Coproduct.} Use deconcatenation on finite bar words.
\item \emph{Compatibility.} Verify that the projection
 $\bar B_X(A_{\le N+1})\to\bar B_X(A_{\le N})$ commutes with
 differential, curvature, coproduct, and the finite-stage twisting
 morphism.
\end{enumerate}
\end{remark}

\subsection{Convergence criteria in practice}

\begin{proposition}[Weight-window convergence; \ClaimStatusProvedHere]\label{prop:practical-convergence}
Let $A=\bigoplus_{h\ge0}A_h$ be a positive-energy chiral algebra with
finite-dimensional weight spaces.  Suppose the chiral $A_\infty$
operations preserve total conformal weight after the reduced-weight
shift used in the bar construction, and the finite weight windows
$A_{\le N}$ lie in the finite-stage bar--cobar regime.  Then
$\widehat{\bar B}(A)=\varprojlim_N\bar B(A_{\le N})$ is well-defined,
and every fixed total weight~$w$ stabilizes at stage~$N=w$.  More
precisely,
\[
\widehat{\bar B}(A)\cong
\prod_{w\ge0}\bar B(A_{\le w})_w,
\qquad
\bar B(A)=\bigoplus_{w\ge0}\bar B(A)_w
\]
with the direct sum sitting as the polynomial subcomplex of the
completed product.  The passage from finite windows to the completed
object is exact only because the weightwise inverse systems are
eventually constant.

Two standard cases follow.
\begin{enumerate}
\item For finitely strongly generated positive-energy algebras, each
 total-weight piece of the bar complex is finite-dimensional.
\item For towers with generators of weights
 $h_1\le h_2\le\cdots\to\infty$, a fixed weight window contains only
 finitely many generators, so the inverse system is eventually
 constant on that window.
\end{enumerate}
\end{proposition}

\begin{proof}
Fix a total weight~$w$. Weight preservation implies that no bar word
of total weight~$w$ can use a generator of weight greater than~$w$.
Thus the projection from the $N$th window to the $w$th window is an
isomorphism on weight~$w$ for all $N\ge w$. The inverse system is
therefore Mittag--Leffler and eventually constant weightwise, and the
completed bar coalgebra is the product of the stabilised weight
pieces. The two stated cases are exactly the two finiteness
conditions ensuring finite-dimensionality before the inverse limit.
\end{proof}

\section{Connection to Costello--Gwilliam}

\subsection{Renormalization and factorization}

\begin{theorem}[Costello--Gwilliam renormalization \cite{CG17}; \ClaimStatusProvedElsewhere]\label{thm:CG-renorm}
\textup{(Costello--Gwilliam, Volume 2)}

For a perturbative quantum field theory, the effective interaction satisfies:
\[I_{\text{eff}} = I_{\text{classical}} + \sum_{n=1}^{\infty} \hbar^n I^{(n)}\]

where $I^{(n)}$ are n-loop quantum corrections.

The bar completion has the same formal inverse-limit shape:
\[
\widehat{\bar B}(A)=\varprojlim_N\bar B(A_{\le N}).
\]

\end{theorem}

\begin{remark}[Physical interpretation]
Under the Feynman diagram identification of Chapter~\ref{ch:v1-feynman},
the completion topology is the algebraic counterpart of a graph-window
expansion: convergence of the inverse system
$\{\bar B(A_{\le N})\}_N$ is exactly the assertion that the formal
graph sum defines a continuous dg coalgebra.
\end{remark}

\section{Higher genus corrections}

\subsection{Genus expansion of completion}

\begin{remark}[Genus-graded completion]\label{rem:genus-completion}
The $I$-adic completion and the genus expansion are two independent filtrations on
the bar complex. At each fixed genus~$g$, the bar complex
$\bar{B}^{(g)}(\mathcal{A})$ receives its own $I$-adic completion:
\[\widehat{\bar{B}}^{(g)}(\mathcal{A}) := \varprojlim_n \bar{B}^{(g)}(\mathcal{A})/I^n.\]
The genus-$g$ contribution involves configuration spaces on Riemann surfaces of
genus~$g$, modular forms, and period integrals over $\mathcal{M}_g$.

The total bar complex collecting all genera is a \emph{separate} object:
\[\bar{B}^{\mathrm{tot}}(\mathcal{A}) := \bigoplus_{g \geq 0} g_s^{2g-2}\, \bar{B}^{(g)}(\mathcal{A}).\]
The convergence of this genus sum requires analytic input beyond the $I$-adic
topology; this is the content of the HS-sewing criterion
(Theorem~\ref{thm:general-hs-sewing}) and not a consequence of the $I$-adic
completion alone. The string-theory interpretation of the genus parameter as a
loop-counting variable $g_s$ is not used in the algebraic
arguments of this chapter.
\end{remark}

%% ====================================================================
%% STABILIZED AND RESONANCE-FILTERED COMPLETION
%% ====================================================================
\section{Stabilized and resonance-filtered completion}
\label{sec:mc4-splitting}

The strong completion-tower theorem
(Theorem~\ref{thm:completed-bar-cobar-strong}) turns the abstract
Mittag--Leffler and continuity inputs of
Proposition~\ref{prop:mc4-reduction-principle} into a finite-window
criterion.  Infinite-generator towers separate into positive towers,
where weight windows stabilize, and resonant towers, where a
finite-dimensional weight-zero sector remains.

\begin{remark}[Positive and resonant towers]
\label{rem:mc4-splitting}
The completed bar-cobar problem for infinite-generator chiral
algebras splits into two independent sub-problems:
\begin{itemize}
\item \emph{Positive towers.}
 When the algebra carries an honest positive weight grading
 (conformal weight, spin, loop degree) and the bar differential
 preserves total weight, the completion problem is reduced by
 weightwise stabilization:
 the weight-$w$ summand of the reduced bar complex depends only on
 the finite truncation $\cA_{\le w}$, and the completed
 counit is a quasi-isomorphism weight by weight.
 This is exactly the content of
 Proposition~\ref{prop:mc4-weight-cutoff} and
 Corollary~\ref{cor:winfty-weight-cutoff}.
 \emph{Positive towers include $\mathcal{W}_{1+\infty}$,
 Drinfeld/RTT positive towers, and standard RTT algebras.}
\item \emph{Finite resonance towers.}
 When the algebra has a finite-dimensional weight-zero
 resonance piece on which higher operations can preserve
 filtration degree~$0$, the stabilization argument does not
 apply directly.  The remaining datum is the finite curved factor
 carrying weight-zero propagation.
 \emph{The Virasoro algebra, $\mathcal{W}_N$ algebras
 ($N \ge 3$), and non-quadratic chiral algebras with central
 terms belong to this regime.}
\end{itemize}
This splitting is structural: the positive sector is controlled by
weight stabilization, while the resonant sector is controlled by a
finite-dimensional curved factor.
\end{remark}

\subsection{Positive towers: stabilized completion}

\begin{theorem}[Stabilized completion for positive towers;
\ClaimStatusProvedHere]
\label{thm:stabilized-completion-positive}
Let $\cA = \prod_{w \ge 0} \cA_w$ be a complete augmented
$A_\infty$-algebra such that:
\begin{enumerate}
\item each $\cA_w$ is finite-dimensional;
\item every $A_\infty$-operation $m_n$ preserves total weight;
\item the curvature $m_0$ has strictly positive weight;
\item each truncation
 $\cA_{\le N} := \prod_{0 \le w \le N} \cA_w$
 lies in the proved finite-type bar-cobar regime.
\end{enumerate}
Then the \emph{stabilized completed reduced bar coalgebra}
\[
\widehat{\bar B}^{\mathrm{st}}(\cA)
:= \prod_{w \ge 0} \bar B(\cA_{\le w})_w
\]
is a well-defined complete curved dg coalgebra, and the
completed counit
\[
\Omega\bigl(\widehat{\bar B}^{\mathrm{st}}(\cA)\bigr)
\longrightarrow \cA
\]
is a quasi-isomorphism in the completed graded category.
\end{theorem}

\begin{proof}
A bar word of total weight~$w$ can only involve inputs of weights
$\le w$. Since the $A_\infty$-operations preserve total weight,
the differential on the weight-$w$ piece sees only $\cA_{\le w}$.
Hence the chain complex $\bar B(\cA)_w$ stabilizes at stage~$w$:
for every $N \ge w$, the natural map
$\bar B(\cA_{\le N})_w \xrightarrow{\sim} \bar B(\cA_{\le w})_w$
is an isomorphism of chain complexes.

The finite-stage counits
$\Omega(\bar B(\cA_{\le N})) \xrightarrow{\sim} \cA_{\le N}$
are quasi-isomorphisms by hypothesis~(4). Restricting to
the weight-$w$ summand (which stabilizes at stage~$w$)
gives a quasi-isomorphism on each weight piece. Taking the
product over~$w$ gives the completed statement.  Products of
weightwise quasi-isomorphisms are exact here because the tower is
eventually constant in each finite-dimensional weight and therefore
strict Mittag--Leffler.
\end{proof}

\begin{remark}[Relation to finite-window criteria]
\label{rem:stabilized-vs-mc4-criteria}
Theorem~\ref{thm:stabilized-completion-positive} is the
positive-tower special case of the strong completion-tower
theorem (Theorem~\ref{thm:completed-bar-cobar-strong}),
under the strong filtration axiom
(Definition~\ref{def:strong-completion-tower}).
The weight-cutoff criterion of
Proposition~\ref{prop:mc4-weight-cutoff} and its application to
the $\mathcal{W}_\infty$ tower in
Corollary~\ref{cor:winfty-weight-cutoff} are concrete instances.
For any target with an honest positive spin, loop, or conformal
grading \textup{(}$\mathcal{W}_{1+\infty}$, Drinfeld/RTT positive
towers, positive RTT algebras\textup{)}, the completed bar-cobar
problem is reduced to finite truncations and strict weightwise
Mittag--Leffler.
\end{remark}

\subsection{Resonance rank}

\begin{definition}[Resonance decomposition and resonance rank]
\label{def:resonance-rank}
Let $\cA$ be a complete filtered $A_\infty$-algebra with a closed
decomposition
\[
\cA \cong R_\cA \;\widehat{\oplus}\; \cA^{>0},
\]
where $\cA^{>0}$ carries a positive weight grading and $R_\cA$ is
the maximal closed subspace on which higher operations can preserve
filtration degree~$0$. Define the \emph{resonance rank}
\[
\rho(\cA) := \dim H^*(R_\cA, m_1^{R_\cA})
\]
when finite.
\begin{itemize}
\item $\rho(\cA) = 0$: purely positive tower;
 stabilization (Theorem~\ref{thm:stabilized-completion-positive})
 solves the completion problem.
\item $0 < \rho(\cA) < \infty$: finite resonance obstruction;
 Theorem~\ref{thm:resonance-filtered-bar-cobar} applies.
\item $\rho(\cA) = \infty$: genuinely wild completion problem.
\end{itemize}
\end{definition}

\begin{remark}[Resonance rank is the true obstruction]
\label{rem:resonance-true-obstruction}
Infinite generation by itself is not the enemy.
The enemy is neutral-energy resonance. The standard
mode-level objection to Virasoro conilpotence (that
coproducts of~$L_0$ spread into infinitely many
$L_k \otimes L_{-k}$ terms) may be a presentation artifact
rather than the intrinsic obstruction of the positive-energy
chiral algebra. The chiral state space of the Virasoro VOA
is positively graded by conformal weight, and each weight space
is finite-dimensional. The right question is not ``is the raw
mode bar construction conilpotent?''\ but rather ``can one
rewrite Virasoro completion on the positive-energy state-space
side, with only a finite resonance slice carrying the
central term and normal-ordering corrections?''
\end{remark}

\begin{lemma}[Exact tensoring with a finite resonance tower; \ClaimStatusProvedHere]
\label{lem:finite-resonance-tensor-exact}
Let
\[
P=\prod_{w\ge0}P_w,\qquad
Q=\prod_{\ell\ge0}Q_\ell
\]
be completed products of complexes over~$\bk$, with every
$P_w$ and $Q_\ell$ finite-dimensional.  The completed tensor product
for the product filtrations is
\[
P\widehat\otimes Q
\cong
\prod_{w,\ell\ge0}(P_w\otimes Q_\ell).
\]
If $f:P\to P'$ and $g:Q\to Q'$ are products of weightwise
quasi-isomorphisms and the corresponding inverse systems are strict
Mittag--Leffler, then
$f\widehat\otimes g$ is a quasi-isomorphism.  No exactness assertion is
made here for arbitrary completed tensor products with an
infinite-dimensional unfiltered factor.
\end{lemma}

\begin{proof}
The displayed tensor product is the inverse limit over finite
rectangles $(w,\ell)\le(N,L)$.  On each rectangle tensoring is exact,
because it is a finite tensor product of complexes of vector spaces.
The transition maps are projections, hence strict epimorphisms.  The
Milnor sequence has vanishing $\varprojlim^1$ by strict
Mittag--Leffler, so cohomology commutes with the limit.  This proves
both the formula and preservation of quasi-isomorphisms.
\end{proof}

\subsection{Resonance-filtered completion}

\begin{theorem}[Resonance-filtered completed bar/cobar;
\ClaimStatusProvedHere]
\label{thm:resonance-filtered-bar-cobar}
Assume $\cA = R_\cA \;\widehat{\oplus}\; \cA^{>0}$ with:
\begin{enumerate}
\item $\cA^{>0}$ satisfies the hypotheses of
 Theorem~\textup{\ref{thm:stabilized-completion-positive}};
\item $R_\cA$ is a finite-dimensional curved $A_\infty$-subalgebra;
\item every mixed higher product with at least one resonance input and
 at least one $\cA^{>0}$-input strictly lowers positive weight;
\item the compatible finite-stage counit cone tower is strict
 Mittag--Leffler.
\end{enumerate}
Define the \emph{resonance-filtered completed bar coalgebra}
\[
\widehat{\bar B}^{\mathrm{res}}(\cA)
:=
\Bigl(\prod_{w \ge 0}
 \bar B\bigl((\cA^{>0})_{\le w}\bigr)_w\Bigr)
\;\widehat{\otimes}\;
\widehat{\bar B}(R_\cA),
\]
where $\widehat{\bar B}(R_\cA)$ is completed in bar length.  Since
$R_\cA$ is finite-dimensional, each resonance bar-length piece is
finite-dimensional.
Then the total bar differential splits as
\[
D = D^+ + D_R + D_{\mathrm{mix}},
\]
where $D^+$ is the stabilized positive-weight differential,
$D_R$ is the bar differential on the resonance-length completion, and
$D_{\mathrm{mix}}$ is
locally nilpotent on every finite positive-weight window
\textup{(}each application strictly lowers positive weight\textup{)}.
Consequently:
\begin{enumerate}
\item $D$ is continuous in the completed topology;
\item $\widehat{\bar B}^{\mathrm{res}}(\cA)$ is a complete
 curved dg coalgebra; and
\item if the finite-stage counits are compatible quasi-isomorphisms,
 then
 \[
 \Omega\bigl(\widehat{\bar B}^{\mathrm{res}}(\cA)\bigr)
 \simeq \cA.
 \]
\end{enumerate}
\end{theorem}

\begin{proof}
Because $D_{\mathrm{mix}}$ strictly lowers positive weight, its
iterates terminate on every finite window: for an element of
positive weight~$w$, the $n$-fold iterate $D_{\mathrm{mix}}^n$
vanishes for $n > w$. Thus the full differential is a finite
perturbation of $D^+ + D_R$ on each window and a continuous
perturbation after completion. The stabilized positive sector is
controlled by Theorem~\ref{thm:stabilized-completion-positive};
the completed tensor product with the resonance sector is exact by
Lemma~\ref{lem:finite-resonance-tensor-exact}; and the mixed term
converges by the weight bound. The filtered homological perturbation
lemma applies on each finite positive-weight and resonance-length
window.  Passing to the completed limit uses the strict
Mittag--Leffler hypothesis on the finite-stage counit cones and
Theorem~\ref{thm:completed-bar-cobar}, giving completed bar-cobar
recovery.
\end{proof}

\begin{construction}[The resonance spectral sequence]
\label{const:resonance-spectral-sequence}
Let $\cA = R_\cA \;\widehat{\oplus}\; \cA^{>0}$ satisfy the
hypotheses of
Theorem~\ref{thm:resonance-filtered-bar-cobar}. Define the
\emph{positive-weight filtration} on the bar complex by
\[
\mathcal{F}_{\le p} \barB(\cA)
:= \bigl\{\text{bar chains whose total weight from $\cA^{>0}$
components is $\le p$}\bigr\}.
\]
Since the mixed differential $D_{\mathrm{mix}}$ strictly
lowers positive weight, this finite increasing filtration is
preserved by $D$, and the associated spectral sequence
$(E_r^{p,q}, d_r)_{r \ge 0}$ converges to
$H^*(\barB(\cA), D)$. The pages are:
\begin{itemize}
\item $E_0$: the associated graded of the bar complex by
 positive weight, with $d_0 = D^+ + D_R$ (the weight-preserving
 part of the differential);
\item $E_1 \cong H^*(\barB(\cA^{>0}), D^+) \otimes
 H^*(\widehat{\barB}(R_\cA), D_R)$, computed by finite
 positive-weight and finite resonance-length windows;
\item $d_1$: the first correction from $D_{\mathrm{mix}}$,
 which lowers positive weight by exactly~$1$;
\item $E_r$ for $r \ge 2$: successive corrections from the mixed
 differential.
\end{itemize}
At total weight~$w$, the spectral sequence degenerates at page
$E_{w+1}$ (since $D_{\mathrm{mix}}$ can lower positive
weight at most $w$ times before reaching~$0$).
In particular, the spectral sequence is
\emph{finitely convergent at every weight level}: each $E_\infty^w$
is computed by a finite number of pages.

\emph{Computational significance.}
The $E_1$ page is computable by finite windows:
$H^*(\barB(\cA^{>0}), D^+)$ is determined by the stabilized
positive-weight theory
(Theorem~\ref{thm:stabilized-completion-positive}), and
$H^*(\widehat{\barB}(R_\cA), D_R)$ is computed by the
bar-length tower of the finite-dimensional resonance algebra. The
differentials $d_r$ for $r \ge 1$ are
explicit: they are determined by the coupling between the
resonance piece and the positive sector through the
mixed $A_\infty$-operations. The spectral sequence provides a
systematic algorithm for computing the bar cohomology of any
finite-resonance algebra.
\end{construction}

\begin{proposition}[Resonance spectral sequence degeneration;
\ClaimStatusProvedHere]
\label{prop:resonance-ss-degeneration}
In the setting of
Construction~\textup{\ref{const:resonance-spectral-sequence}},
if the mixed operations
$\cA^{>0} \otimes R_\cA \to \cA^{>0}$ lower positive weight
by at least~$\delta \ge 1$, then the spectral sequence
degenerates at page $E_{\lfloor w/\delta \rfloor + 1}$ for
elements of total weight~$w$.

In particular, if $\delta = 1$ the spectral sequence for
weight-$w$ elements has at most $w+1$ pages, and if
$\delta \ge w$ it degenerates at $E_2$
\textup{(}the first correction suffices\textup{)}.
\end{proposition}

\begin{proof}
Each non-zero application of $d_r$ lowers positive weight by at
least~$r \cdot \delta$. For an element of positive weight~$w$,
no term can be lowered more than~$w$, so $d_r = 0$ for
$r > w/\delta$. Hence $E_{\lfloor w/\delta \rfloor + 1}
= E_\infty$ for weight-$w$ elements.
\end{proof}

\subsection{Resonance rank computations for standard families}

\begin{proposition}[Resonance ranks of the standard families;
\ClaimStatusProvedHere]
\label{prop:resonance-ranks-standard}
For the standard chiral algebra families, the resonance rank
and resonance decomposition are:
\begin{center}
\renewcommand{\arraystretch}{1.3}
\begin{tabular}{lccl}
\textbf{Family} & $\rho(\cA)$ & $\dim R_\cA$ & \textbf{Notes} \\
\hline
Heisenberg $\mathcal{H}_k$ & $0$ & $0$
 & Quadratic; no completion needed \\
Free fermion $\mathrm{Fer}$ & $0$ & $0$
 & Quadratic; no completion needed \\
$\beta\gamma$ system & $0$ & $0$
 & Quadratic; no completion needed \\
$\widehat{\mathfrak{g}}_k$ \textup{(}affine\textup{)} & $0$ & $0$
 & Quadratic OPE; PBW-Koszul \\
$\mathrm{Vir}_c$ \textup{(}Virasoro\textup{)} & $1$ & $1$
 & $R = \bk \cdot m_0$
 \textup{(}curvature from quartic pole\textup{)} \\
$\mathcal{W}_3$ & $1$ & $1$
 & $R = \bk \cdot m_0$
 \textup{(}curvature from sixth-order pole\textup{)} \\
$\mathcal{W}_N$ \textup{(}$N \ge 3$\textup{)} & $1$ & $1$
 & Curvature $m_0 \propto c$;
 single resonance direction \\
$\mathcal{W}_{1+\infty}$ & $0$ & $0$
 & Positive conformal-weight tower;
 stabilization applies \\
$Y_{\mathrm{RTT}}(\mathfrak{g})$ \textup{(}Drinfeld/RTT tower\textup{)} & $0$ & $0$
 & Positive loop-weight tower;
 stabilization applies
\end{tabular}
\end{center}
\end{proposition}

\begin{proof}
\emph{Quadratic families} ($\mathcal{H}_k$, $\mathrm{Fer}$,
$\beta\gamma$, $\widehat{\mathfrak{g}}_k$):
The OPE is at most quadratic, and the bar complex is
finite-dimensional at each weight. The strong completion-tower
axiom (Definition~\ref{def:strong-completion-tower}) is
satisfied for the PBW filtration; no resonance piece is needed.

\emph{Virasoro and $\mathcal{W}_N$} ($N \ge 3$):
In the curved $A_\infty$-decomposition of the chiral algebra
structure, the central charge contributes a curvature element
$m_0 = \kappa(\cA) \cdot \eta \in F^0\cA$, where $\eta$ is
the vacuum covector.
The resonance piece $R_\cA = \bk \cdot m_0$ is
one-dimensional, with trivial differential $d_R(m_0) = 0$
(the curvature is closed by the $A_\infty$-identity
$m_1(m_0) + m_2(m_0, m_0) = 0$, and for a single generator
this reduces to $m_1(m_0) = 0$ plus higher-order terms that
vanish by dimension). Hence $\rho(\cA)
= \dim H^*(R_\cA, m_1^R) = 1$. The finite-dimensionality
is direct from the positive-energy state-space grading:
$R_\cA=H_0$, each $V_h$ is finite-dimensional, and for Virasoro
$\dim H_0 \leq \dim V_0 = 1$.

\emph{Positive towers}
($\mathcal{W}_{1+\infty}$ and Drinfeld/RTT positive towers):
The tower structure provides a positive conformal-weight (resp.\
RTT-weight) grading. The curvature at each finite stage is
absorbed into the finite-type bar-cobar regime, and the
stabilized completion of
Theorem~\ref{thm:stabilized-completion-positive} applies.
The infinite-generator limit has $R_\cA = 0$ because the
curvature at each stage is handled by the finite truncation.
This row concerns ordinary Drinfeld/RTT towers.  It is not the K3
Hall--Drinfeld or BKM target object, whose completion data live in the
separate Vol.~III comparison lane.
\end{proof}

\begin{corollary}[Virasoro resonance spectral sequence;
\ClaimStatusProvedHere]
\label{cor:virasoro-resonance-ss}
Assume $\mathrm{Vir}_c$ admits a resonance splitting with
$R_{\mathrm{Vir}}=\bk\cdot m_0$, that the resonance bar factor has
$H^*(\widehat{\barB}(R_{\mathrm{Vir}}),D_R)\cong\bk$, and that the
finite-stage counit cones are strict Mittag--Leffler.  Then the
resonance spectral sequence of
Construction~\textup{\ref{const:resonance-spectral-sequence}}
has:
\begin{itemize}
\item $E_1 \cong H^*(\barB(\mathrm{Vir}_c^{>0}))$
 \textup{(}since the resonance bar cohomology is scalar by
 hypothesis\textup{)};
\item if a separate Verdier--Koszul finite-stage computation identifies
 $H^*(\barB(\mathrm{Vir}_c^{>0}))$ with the same-family shadow
 $\mathrm{Vir}_{26-c}$, then $\mathrm{Vir}_{26-c}$ is this
 $E_1$ page;
\item the differentials $d_r$ for $r \ge 1$ are the corrections
 from the central charge coupling.
\end{itemize}
Thus the same-family shadow is a first-page identification only under
the stated Verdier--Koszul input; it is not the completed dual by
itself.
\end{corollary}

\begin{proof}
The general resonance spectral sequence gives
\[
E_1\cong
H^*(\barB(\mathrm{Vir}_c^{>0}),D^+)\otimes
H^*(\widehat{\barB}(R_{\mathrm{Vir}}),D_R).
\]
The scalar-resonance hypothesis reduces the second factor to~$\bk$.
The optional identification with $\mathrm{Vir}_{26-c}$ is not a
consequence of completion; it is exactly the additional
Verdier--Koszul computation named in the statement.  The higher
differentials are induced by $D_{\mathrm{mix}}$, hence by the central
charge coupling.
\end{proof}

\begin{remark}[The Virasoro resonance model]
\label{rem:virasoro-resonance-model}
The Virasoro state-space bar complex carries a positive-energy
filtration.  In the transferred models considered here the central
term gives a one-dimensional resonance candidate
$R_{\mathrm{Vir}} = \mathbb{k} \cdot m_0$.  The resonance construction
uses four mathematical inputs:
\begin{enumerate}
\item \emph{Finite-dimensionality of the resonance slice}: this is a
 hypothesis of the transferred model in
 Theorem~\ref{thm:platonic-completion}
 \textup{(}$R_\cA = H_0$, $\dim H_0 \le \dim V_0 < \infty$\textup{)}.
 For Virasoro: $\dim V_0 = 1$ (just the vacuum), so
 $\dim R_{\mathrm{Vir}} \le 1$.
 Proposition~\ref{prop:resonance-ranks-standard} computes
 $\dim R_{\mathrm{Vir}} = 1$.
\item \emph{Strict positive-weight drop of mixed
 operations}: this is the reduced-weight hypothesis in
 Theorem~\ref{thm:platonic-completion}; once imposed, the mixed
 differential strictly lowers positive conformal weight.
\item \emph{Compatibility with $\mathcal{W}_N$
 approximants}: this holds when the DS-reduction functor preserves the
 positive-energy grading (the BRST differential commutes
 with~$L_0$), so that the resonance decomposition is compatible
 with the $N \to \infty$ limit.
\item \emph{Exactness through functors}: resonance-filtered completion
 commutes with a weight-preserving functor only when that functor is
 exact on the finite windows and preserves the strict
 Mittag--Leffler cone tower.
\end{enumerate}
Under these inputs, the same-family shadow $\mathrm{Vir}_{26-c}$ can
occur only as the first-page Verdier--Koszul shadow named in
Corollary~\ref{cor:virasoro-resonance-ss}.  The completed dual is
computed from the full object
$\widehat{\bar B}^{\mathrm{res}}(\mathrm{Vir}_c)$, with the higher
pages carrying the central-charge corrections.
\end{remark}

\subsection{The Resonance completion theorem}

\begin{theorem}[Resonance completion; \ClaimStatusProvedHere]
\label{thm:platonic-completion}%
\index{platonic completion theorem|textbf}%
\index{resonance rank!finiteness|textbf}%
Let $\cA$ be a separated complete chiral $A_\infty$-algebra obtained
from a positive-energy chiral algebra
$\mathcal V=\prod_{h\ge0}V_h$ by a weight-preserving strong
deformation retract, with every $V_h$ finite-dimensional.  Assume:
\begin{enumerate}
\item the transferred operations satisfy the residue-weight identity
 \[
 h\bigl(\tilde m_k(x_1,\ldots,x_k)\bigr)
 =
 \sum_i h(x_i)-(k-1)
 \]
 whenever the right-hand side is non-negative, and the operation
 vanishes otherwise;
\item every finite reduced-weight truncation lies in the finite-stage
 bar-cobar regime; and
\item the finite-stage counit cone tower is strict Mittag--Leffler.
\end{enumerate}
Then $\cA$ admits the closed decomposition
\[
\cA \cong R_\cA \;\widehat{\oplus}\; \cA^+_\cA
\]
where $\cA^+_\cA$ satisfies stabilized completion
\textup{(}Theorem~\textup{\ref{thm:stabilized-completion-positive})}
and $R_\cA$ is finite-dimensional. The completed reduced bar
coalgebra is then
\[
\widehat{\bar B}(\cA)
\simeq
\widehat{\bar B}^{\mathrm{st}}(\cA^+_\cA)
\;\widehat{\otimes}\;
\widehat{\bar B}(R_\cA),
\]
and the completed cobar recovers~$\cA$.

For $\mathcal{W}_{1+\infty}$, Drinfeld/RTT positive towers, and
positive RTT algebras satisfying these hypotheses, one has
$R_\cA = 0$.
For Virasoro and non-quadratic $\mathcal{W}$-families, one has
$0 < \dim R_\cA < \infty$ when the stated resonance model has been
verified.  Any identification of the same-family shadow
$\mathrm{Vir}_{26-c}$ is a separate Verdier--Koszul computation of
the first page, not a consequence of the decomposition alone.
\end{theorem}

\begin{proof}
Fix a positive-energy chiral algebra $\mathcal V$ whose
weight-preserving transferred $A_\infty$-model is~$\cA$.
Write
\[
\mathcal V = \prod_{h \ge 0} V_h,
\qquad
H_h := H^*(V_h, m_1|_{V_h}),
\qquad
H := \prod_{h \ge 0} H_h.
\]
Choose a strong deformation retract on each weight space and
assemble them into a weight-preserving SDR
\[
(\mathcal V, H, p, \iota, h).
\]
By the homotopy transfer theorem
\textup{(}Theorem~\ref{thm:htt}\textup{)}, $H$ inherits
transferred operations $\tilde m_k$, and by construction
$\cA = (H, \{\tilde m_k\})$.

\emph{Step~1: Reduced weight is preserved on the positive sector.}
For $x \in H_h$ with $h > 0$, define the reduced weight
\[
\nu(x) := h - 1.
\]
The residue-weight hypothesis gives, for homogeneous
$x_i\in H_{h_i}$,
\[
h\bigl(\tilde m_k(x_1, \ldots, x_k)\bigr)
=
\sum_{i=1}^k h_i - (k-1)
\qquad (k \ge 2).
\]
Hence
\[
\nu\bigl(\tilde m_k(x_1, \ldots, x_k)\bigr)
=
\sum_{i=1}^k \nu(x_i)
\qquad (k \ge 2).
\]
Since $H$ is cohomology, the transferred differential vanishes:
$\tilde m_1 = 0$. The transferred curvature $\tilde m_0$ has
conformal weight~$0$.

Set
\[
R_\cA := H_0,
\qquad
\cA^+_\cA := \prod_{h > 0} H_h.
\]
Then $\cA = R_\cA \;\widehat{\oplus}\; \cA^+_\cA$. The formula
above shows that $\cA^+_\cA$ is closed under all
$\tilde m_k$ with $k \ge 1$, because if every input has
conformal weight at least~$1$, then the output has conformal
weight at least~$1$ as well. On $\cA^+_\cA$, the transferred
curvature is zero, so the curvature clause in
Theorem~\ref{thm:stabilized-completion-positive} is vacuous.
Each reduced-weight truncation
\[
(\cA^+_\cA)_{\le N}^{\nu}
:=
\prod_{0 \le \nu \le N} (\cA^+_\cA)_\nu
=
\prod_{1 \le h \le N+1} H_h
\]
is finite-dimensional and lies in the finite-stage bar-cobar regime
by hypothesis. Therefore
Theorem~\ref{thm:stabilized-completion-positive} applies to
$\cA^+_\cA$ with respect to the reduced weight~$\nu$.

\emph{Step~2: The resonance piece is finite-dimensional and
curved.}
By positive energy, each $V_h$ is finite-dimensional. Therefore
each $H_h$ is finite-dimensional, in particular
$R_\cA = H_0$ is finite-dimensional. If all inputs lie in $H_0$
and $k \ge 2$, then the conformal-weight formula from Step~1 gives
\[
h\bigl(\tilde m_k(x_1, \ldots, x_k)\bigr) = -(k-1) < 0,
\]
which is impossible in a positive-energy theory. Hence
\[
\tilde m_k|_{R_\cA^{\otimes k}} = 0
\qquad (k \ge 2).
\]
Since $\tilde m_1 = 0$ and $\tilde m_0 \in R_\cA$, the resonance
piece is a finite-dimensional curved $A_\infty$-subalgebra.
Its resonance rank is therefore
\[
\rho(\cA)
=
\dim H^*(R_\cA, \tilde m_1)
=
\dim R_\cA
<
\infty.
\]

\emph{Step~3: The mixed differential is topologically nilpotent.}
Consider the completed tensor product
\[
\widehat{\bar B}^{\mathrm{res}}(\cA)
=
\widehat{\bar B}^{\mathrm{st}}(\cA^+_\cA)
\;\widehat{\otimes}\;
\widehat{\bar B}(R_\cA).
\]
For a bar word, define its \emph{positive conformal weight}
$w_+$ to be the sum of the conformal weights of the
$\cA^+_\cA$-entries. The differential splits as
$D = D^+ + D_R + D_{\mathrm{mix}}$, where $D^+$ uses only
operations on~$\cA^+_\cA$, $D_R$ uses only operations on
$R_\cA$, and $D_{\mathrm{mix}}$ contains the remaining terms.

If a mixed operation involves $r \ge 1$ positive inputs of
weights $h_1, \ldots, h_r$ and $s \ge 1$ resonance inputs, then
\[
h_{\mathrm{out}}
=
\sum_{i=1}^r h_i - (r+s-1)
\le
\sum_{i=1}^r h_i - 1.
\]
So any mixed operation with at least one resonance input lowers
positive conformal weight by at least~$1$. If all inputs are
positive, then $h_{\mathrm{out}} \ge 1$, so the output remains in
$\cA^+_\cA$ and the term belongs to~$D^+$, not to
$D_{\mathrm{mix}}$. Consequently
$D_{\mathrm{mix}}$ lowers $w_+$ by at least~$1$ on every term to
which it contributes.

Therefore, for a bar chain of positive conformal weight~$w$,
\[
D_{\mathrm{mix}}^n = 0
\qquad\text{for } n > w.
\]
This is the required topological nilpotence. The perturbation
argument used in the proof of
Theorem~\ref{thm:resonance-filtered-bar-cobar} depends only on
that nilpotence, so it applies here with the finite increasing
filtration
\[
\mathcal F_{\le p}
:=
\{\text{bar chains of positive conformal weight } \le p\}.
\]
Combining this with Step~1 gives
\[
\widehat{\bar B}(\cA)
\simeq
\widehat{\bar B}^{\mathrm{st}}(\cA^+_\cA)
\;\widehat{\otimes}\;
\widehat{\bar B}(R_\cA),
\]
and the completed cobar recovers~$\cA$.

\emph{Step~4: Comparison with the original positive-energy model.}
The SDR maps extend to $A_\infty$-quasi-isomorphisms by
Theorem~\ref{thm:htt}. By
Proposition~\ref{prop:transfer-bar}, the induced maps on bar
complexes are quasi-isomorphisms:
\[
\bar B(\mathcal V) \xrightarrow{\sim} \bar B(\cA),
\qquad
\bar B(\cA) \xrightarrow{\sim} \bar B(\mathcal V).
\]
Because the SDR data preserve conformal weight, these maps are
weightwise quasi-isomorphisms.  Passing to the completed products over
weight and resonance-length pieces preserves quasi-isomorphisms by
Lemma~\ref{lem:finite-resonance-tensor-exact} and the strict
Mittag--Leffler cone hypothesis. Hence the resonance-filtered
completed bar coalgebra constructed above computes the same completed
bar-cobar type as the original positive-energy chiral
algebra~$\mathcal V$.

The specific resonance ranks for the standard families are
computed in Proposition~\ref{prop:resonance-ranks-standard}:
$\rho = 0$ for quadratic and positive-tower families, and
$\rho = 1$ for Virasoro and $\mathcal{W}_N$
\textup{(}the single resonance direction being the
curvature $m_0$\textup{)}.
\end{proof}

\begin{remark}[The mode algebra vs.\ state-space formulation]
\label{rem:mode-vs-state-completion}
The proof uses the \emph{state-space} formulation of the
chiral algebra (positive-energy grading on the VOA state space),
not the \emph{mode algebra} formulation (generators $L_n$ for
all~$n \in \mathbb{Z}$). In the mode algebra, the coproduct
of~$L_0$ involves infinitely many terms
$L_k \otimes L_{-k}$, and the raw mode bar construction has
infinite-dimensional weight-$0$ part. The positive-energy
state-space formulation avoids this: each weight space is
finite-dimensional by axiom, and the SDR is constructed weight by
weight. The passage from modes to states is the HTT
(Theorem~\ref{thm:chiral-htt}), and the weight compatibility
(Step~2 above) ensures that the resonance piece remains
finite-dimensional after transfer.
\end{remark}

\begin{construction}[Completion towers as strict inverse limits]
\label{constr:completion-mc4}
The strong completion-tower theorem
(Theorem~\ref{thm:completed-bar-cobar-strong}) uses the degree cutoff
lemma
(Lemma~\ref{lem:degree-cutoff}): the strong filtration axiom
\[
\mu_r(F^{i_1}, \ldots, F^{i_r})
\subset F^{i_1 + \cdots + i_r}
\]
ensures continuity and strict Mittag--Leffler for the finite-window
tower, yielding the completed bar-cobar homotopy equivalence on
$\operatorname{CompCl}(\mathcal{F}_{\mathrm{ft}})$. The
resonance rank $\rho(\cA)$
(Definition~\ref{def:resonance-rank}) classifies completion
difficulty.
\begin{enumerate}[label=\textup{(\roman*)}]
\item \emph{Positive towers, $\rho = 0$.}
 Weight stabilization
 (Theorem~\ref{thm:stabilized-completion-positive})
 reduces the completed bar-cobar problem to a finite
 verification at each weight: the transition maps
 $\widehat{\bar B}_{k+1}/F^{n+1}
 \twoheadrightarrow
 \widehat{\bar B}_k/F^n$
 stabilize for $n$ exceeding a bound depending only on
 the weight grading, so
 $\varprojlim^1 = 0$ and the inverse limit is exact.
 Examples: $\mathcal{W}_{1+\infty}$, Drinfeld/RTT positive
 towers, RTT algebras.
\item \emph{Finite resonance towers, $\rho > 0$.}
 The resonance-filtered bar-cobar
 (Theorem~\ref{thm:resonance-filtered-bar-cobar})
 decomposes the problem into a finite-dimensional
 resonance piece $R_\cA$ and a positive complement:
 $\widehat{\bar B}(\cA)
 \simeq
 \widehat{\bar B}^{\mathrm{st}}(\cA^+)
 \,\widehat{\otimes}\,
 \widehat{\bar B}(R_\cA)$.
 The remaining completion difficulty is concentrated in the
 bar-length tower of the finite-dimensional resonance algebra.
 Examples: $\mathrm{Vir}_c$ ($\rho = 1$),
 non-quadratic $\mathcal{W}_N$ ($\rho < \infty$,
 Theorem~\ref{thm:platonic-completion}).
\item The coefficient-stability criterion
 (Theorem~\ref{thm:coefficient-stability-criterion})
 reduces convergence of the completed bar-cobar to
 finite matrix stabilization at each weight window,
 and the uniform PBW bridge
 (Theorem~\ref{thm:uniform-pbw-bridge}) connects PBW concentration
 to completed bar-cobar inversion on the stated finite-window locus.
\end{enumerate}
\end{construction}

\begin{remark}[Coideal descent and resonance factorisation]
\label{rem:mc4-coideal-descent}
The splitting of Remark~\ref{rem:mc4-splitting} imposes the following
mathematical order on quotient constructions:
\begin{enumerate}
\item Positive towers are governed by
 Theorem~\ref{thm:stabilized-completion-positive}; the completed
 counit is checked weight by weight.
\item For Drinfeld/RTT and $\mathcal{W}_{1+\infty}$-type targets, the
 unshifted completed envelope precedes shifted quotients.  The
 kernel $\to$ completed envelope $\to$ boundary quotient order is
 forced by the coideal structure
 \textup{(}Appendix~\textup{\ref{chap:shifted-rtt-orthogonal-coideals}}\textup{)}.
\item Orthogonal coideals and shifted quotients descend through
 completion only after their kernels are closed for the
 finite-window topology.  Shifted BFN-type data belong to this
 quotient/coideal descent problem.
\item The Virasoro state-space bar complex carries the
 positive-energy filtration of
 Remark~\ref{rem:virasoro-resonance-model}; its completed bar object
 is controlled by the finite-dimensional resonance factor
 $R_{\mathrm{Vir}}$.
\item PBW concentration supplies the finite windows, completed
 bar--cobar supplies the inverse limit, and the BV/BRST dictionary is a
 later interpretation of that completed object.
\end{enumerate}
\end{remark}

\section{Summary}

\begin{remark}[Completed finite-window package]\label{rem:nilpotent-main}
The completed bar--cobar package consists of:
\begin{enumerate}
\item the finite-stage convergence and geometry of the completed bar
 tower
 \textup{(}Theorem~\ref{thm:completion-convergence},
 Proposition~\ref{prop:geom-conilpotent}\textup{)};
\item the inverse-limit criteria used by the standard
 $\mathcal{W}_\infty$ and Drinfeld/RTT positive towers;
\item the compatibility of these completed models with the
 Beilinson--Drinfeld filtered framework
 \textup{(}Theorem~\ref{thm:BD-chiral-homology}\textup{)};
\item the splitting into positive and finite-resonance regimes
 \textup{(}Remark~\ref{rem:mc4-splitting}\textup{)};
\item stabilized completion for positive towers
 \textup{(}Theorem~\ref{thm:stabilized-completion-positive}\textup{)};
\item the resonance rank invariant
 \textup{(}Definition~\ref{def:resonance-rank}\textup{)}
 and resonance-filtered completed bar-cobar
 \textup{(}Theorem~\ref{thm:resonance-filtered-bar-cobar}\textup{)};
\item the Resonance completion theorem
 \textup{(}Theorem~\ref{thm:platonic-completion}\textup{)},
 which proves, under reduced-weight and strict Mittag--Leffler
 hypotheses, that the completed bar-cobar type is represented by a
 finite-resonance transferred model.
\end{enumerate}
The Resonance completion theorem
\textup{(}Theorem~\ref{thm:platonic-completion}\textup{)}
is the sharp finite-window form for positive-energy transferred models
satisfying its hypotheses.  The remaining algebraic computation is the
resonance spectral sequence for $\mathrm{Vir}_c$ and the
non-quadratic $\mathcal W_N$ algebras, together with the
Verdier--Koszul identification of the H-level target from its higher
pages.
\end{remark}

\section{Synthesis}
\label{sec:nc-supplement}

\begin{remark}[Siegel-elliptic comparison kernels and Hall completion]
\label{rem:nc-1}
For the K3 Hall comparison one may posit a Siegel-elliptic dynamical
kernel
\[
 R_{\mathrm{Sieg,dyn}}(z_1, z_2;\,\tau, w):\
 V(z_1)\otimes V(z_2) \to V(z_2)\otimes V(z_1)
\]
with spectral parameters $(z_1,z_2)$ on the configuration side and
dynamical parameters $(\tau,w)$ on $\mathbb H_2$.  This chapter does
not construct such a kernel and does not identify it with an ordinary
Drinfeld Yangian $R$-matrix.  The nilpotent completion of the chiral
Hall algebra
$\mathcal Y^{\mathrm{Hall}}_\hbar(\mathrm{CoHA}_{K3\times E})$
supplies only the completed ambient in which a Vol.~III
Hall--Drinfeld target $\mathbf H_{\Delta_5}^{\mathrm{tgt}}$ can be
tested after the target-side comparison maps have been constructed.
Recognition as the compact source $\mathbf H_{\Delta_5}$ is a
Vol.~III source-gate statement, not a construction in this chapter
\cite{PasolZagier2013,SchiffmannVasserot2012,Davison2017}. Cross-ref
Vol.~III Chapter~\ref{V3-ch:k3-chiral-bialgebra-platonic}.
\end{remark}

\begin{remark}[K3 comparison constants do not bound nilpotence]
\label{rem:nc-2}
The identity
\[
\hbar^2 K^{\kappa_{\mathrm{ch}}}=-1,\qquad
K^{\kappa_{\mathrm{ch}}}=8,\qquad
\hbar^2=-1/8
\]
is a comparison normalisation for the K3 target lane.  It is not a
nilpotence bound and it does not imply that the completed bar tower
truncates at order~$8$.  The Vol.~I completion remains the inverse
limit over finite windows.

The K3 constants occupy different registers:
\[
K^{\kappa_{\mathrm{ch}}}=8=2c_+,\quad c_+=4,\quad
\kappa_{\mathrm{ch}}^{\mathrm{Heis}}(K3\times E)=3,\quad
\kappa_{\mathrm{BKM}}(\Delta_5)=5,
\]
while $\kappa_{\mathrm{cat}}(K3\times E)=0$,
$\kappa_{\mathrm{fiber}}(K3)=24$, and $\chi(\mathcal O_{K3})=2$ are
separate invariants.  The numerical equality $5=3+2$ is the
$N=1$ coincidence recorded in the concordance, not a product formula.
The automorphic forms $\Delta_5$ and $\Phi_{10}^{\mathrm{un}}$ are
target comparison outputs for the Vol.~III BKM/Hall lane, not input
data proving a truncation of nilpotent completion
\cite{Lusztig1990,GritsenkoNikulin1998}. Cross-ref
Chapter~\ref{V3-ch:k3-chiral-bialgebra-platonic}.
\end{remark}

\section{Higher-coherence obstruction tower}
\label{sec:nc-supplement-ii}

\begin{remark}[Pentagon target cocycle at $\hbar^3, \hbar^4, \hbar^5$]
\label{rem:nc-phi3-4-5}
The pentagon target equation for the candidate Siegel--Borcherds
scalar cocycle
$\widetilde{\Phi}^{\mathrm{Sieg\text{-}Bor}}_{\mathrm{Sp}_4}[\Phi_{10}^{\mathrm{un}}/\eta^{24}]$
closes at $\hbar^{\le 3}$ with the explicit three-coboundary
\[
\phi^{(3)} \;=\; \zeta(3)\,c_{\mathrm{symm}}
 \;+\;(25/3)\,c_{\mathrm{timelike}}
 \;+\;(\Phi_{10}^{\mathrm{un}}/\eta^{24})\cdot c_{\Phi_{10}^{\mathrm{un}}},
\qquad c_{\Phi_{10}^{\mathrm{un}}}^{\mathrm{leading}} = -2,
\]
the leading Fourier coefficient $-2$ being the
Gritsenko--Nikulin 1997 Theorem~2.1 singular-theta lift coefficient
at $(q_\rho^1 q_\tau^1 y^0)$, with Eguchi--Ooguri--Tachikawa 2011
normalisation $c_{\phi_{0,1}^{K3}}(-1) = 2$ entering with the
singular-theta involution sign
\cite{GritsenkoNikulin1997,EguchiOoguriTachikawa2011}.
The obstruction tower extends to $\hbar^4,\hbar^5$:
\[
\phi^{(4)} \;=\; \tfrac{1}{24}\,\zeta(3)\,c_4^{(1)}
 \;+\;\tfrac{1}{6}\,(\Phi_{10}^{\mathrm{un}}/\eta^{24})^2\,c_4^{(K3)},
\qquad
\phi^{(5)} \;=\; \tfrac{1}{240}\,\zeta(5)\,c_5^{(1)}
 \;+\;\tfrac{1}{120}\,(\Phi_{10}^{\mathrm{un}})^{5/2}/\eta^{60}\,c_5^{(K3)}.
\]
At weight~$4$ no new multiple-zeta irreducibles appear (Brown 2011
``Mixed Tate motives over $\mathbb Z$''); the rational coefficient
$1/24$ is fixed by the Etingof--Kazhdan 2000 super-Drinfeld associator
$(2\pi i)^4/4!$ normalisation of the 4-leg Knizhnik--Zamolodchikov
iterated integral. At weight~$5$ the new irreducible is $\zeta(5)$;
the would-be $\zeta(3)^2$ leg vanishes at $\hbar^5$ because
$\zeta(3)^2$ is a weight-$6$ object (it re-enters at $\hbar^6$,
reducible via the Zagier double-shuffle relations).
\cite{Brown2011,EtingofKazhdan2000,EnriquezGomezGonzalezMaassarani2022,
 GritsenkoNikulin1998,Drinfeld1990,MarklShniderStasheff2002}.
An all-order associator on the compact Hall realisation additionally
requires the Hall source and the all-order pentagon datum.
\end{remark}

\begin{remark}[$A_\infty$-quasi-Hopf: the tower does not truncate]
\label{rem:nc-Ainfty-open}
The conclusion is structural and belongs to the target lane.  For BKM
target inputs whose imaginary-root cone is infinite-dimensional
\textup{(}Borcherds 1992; Gritsenko--Nikulin 1998\textup{)}, the
Vol.~III comparison may have higher-coherence contributions in
unbounded $\hbar$-orders.  Under the Vol.~III imaginary-root counting
hypothesis the target-side dimensions have asymptotic form
\[
\dim\,c_n^{(K3)} \;\sim\; n^{-27/4}\,\exp\!\bigl(4\pi\sqrt{n}\bigr)
\]
by Hardy--Ramanujan asymptotics applied to the Gritsenko--Nikulin~1998
Section~4 imaginary-root counting formula
\cite{GritsenkoNikulin1998,Borcherds1992,HardyRamanujan1918}.
This is not a Vol.~I proof that every compact source has a non-zero
class in every order.  It says that no finite nilpotence bound follows
from the K3 comparison constants alone.  If a compact source is
recognised by the Vol.~III
$\mathbf H_{\Delta_5}^{\mathrm{tgt}}$ comparator and has non-zero
classes in unbounded orders, it carries a genuine
$A_\infty$-quasi-Hopf structure rather than only a quasi-Hopf structure
in Drinfeld's original sense.  The
controlling universal object is the chain-level Markl--Shnider--Stasheff
complex $\mathrm{MSS}^\bullet(\widehat{\mathfrak g}^{\mathrm{super}}_{\Delta_5})$
over the \emph{pro-}Lie-bialgebra
$t^{\mathrm{Sieg},\mathrm{super}}_{2,[2]} \oplus \mathfrak n_+^{\mathrm{imag}}$
(the super-Manin pair of the K3 BKM).  The nilpotent completion of the
present chapter supplies one pro-conilpotent weight-completed ambient
for this comparison; it does not recognise the compact Hall source.
Ambient-qualifier
discipline applies: the chain-level statement holds on the
pro-object ambient; the $(\infty,1)$-categorical statement holds in
the category of presentable $A_\infty$-Hopf $\infty$-operads
(each lane inscribes its own ambient; the two are distinct
theorems, both load-bearing). Cross-ref Vol.~III
Chapter~\ref{V3-ch:k3-chiral-bialgebra-platonic}.
\end{remark}

\section{Deep motivic coherences $\phi^{(8,9,10)}$}
\label{sec:nc-supplement-iii}

\begin{remark}[Brown's motivic MZV dimensions through weight~$12$;
\ClaimStatusProvedElsewhere]
\label{rem:nc-brown-deligne}
Brown's Hoffman basis theorem gives
\[
 \sum_{n\geq0}d_nt^n=\frac{1}{1-t^2-t^3},
 \qquad
 (d_0,d_1,d_2)=(1,0,1).
\]
Hence
\[
d_3=1,\quad d_4=1,\quad d_5=2,\quad d_6=2,\quad
d_7=3,\quad d_8=4,\quad d_9=5,\quad d_{10}=7,\quad
d_{11}=9,\quad d_{12}=12.
\]
The basis in weight~\(n\) is indexed by compositions of \(n\) with
parts in \(\{2,3\}\).  These dimensions describe the motivic period
algebra.  Their image in a pentagon or cyclic deformation complex is
defined by a chosen associator and a word-to-cochain map.
\cite{Brown2011,Hoffman1997}.
\end{remark}

\begin{remark}[Weight-eight period coordinates]
\label{rem:nc-phi8}
The four Hoffman words in weight~$8$ are
\[
 (2,3,3),\quad(3,2,3),\quad(3,3,2),\quad(2,2,2,2).
\]
Thus \(d_8=4\).  The scalar Borcherds section
\[
 \left(\frac{\Delta_5}{\eta^{12}}\right)^8
\]
has raw modular weight \(-8\) and trivial
\(\Delta_5\)-character.  Given lower cochains
\(\Theta_{<8}\), an order-$8$ cyclic operation is a cochain
\(\psi_8\) satisfying
\[
 D\psi_8
 =
 -\frac12\!\sum_{\substack{i+j=8\\i,j\geq3}}
 [\psi_i,\psi_j],
 \qquad
 \operatorname{Rot}_8(\psi_8)=0.
\]
A word-to-cochain map and a Borcherds comparison homotopy place the
four period coordinates and the scalar section in this equation.
\end{remark}

\begin{remark}[Weight-nine period coordinates]
\label{rem:nc-phi9}
The five Hoffman words in weight~$9$ are
\[
 (3,3,3)
 \quad\text{and the four permutations of}\quad
 (2,2,2,3).
\]
Thus \(d_9=5\).  The scalar section
\((\Delta_5/\eta^{12})^9\) has raw modular weight \(-9\) and carries
the character \(\nu_{\Delta_5}\).  An order-$9$ cyclic operation is
a cochain \(\psi_9\) solving
\[
 D\psi_9
 =
 -\frac12\!\sum_{\substack{i+j=9\\i,j\geq3}}
 [\psi_i,\psi_j],
 \qquad
 \operatorname{Rot}_9(\psi_9)=0.
\]
The exact period count and the cyclic solution are related by the
chosen word-to-cochain map.
\end{remark}

\begin{remark}[Weight-ten period coordinates]
\label{rem:nc-phi10}
The seven Hoffman words in weight~$10$ are the six permutations of
\((2,2,3,3)\) together with \((2,2,2,2,2)\).  Hence \(d_{10}=7\).
The scalar section
\((\Delta_5/\eta^{12})^{10}\) has raw modular weight \(-10\) and
trivial character.  The order-$10$ obstruction and solution have the
same form:
\[
 R_{10}:=
 \frac12\!\sum_{\substack{i+j=10\\i,j\geq3}}
 [\psi_i,\psi_j],
 \qquad
 D\psi_{10}=-R_{10},
 \qquad
 \operatorname{Rot}_{10}(\psi_{10})=0.
\]
The Hoffman basis indexes period coordinates; the represented
cochain \(\psi_{10}\) is supplied by the comparison map and the
chosen primitive of \(R_{10}\).
\end{remark}

\begin{remark}[Two asymptotic sequences]
\label{rem:nc-asymptotic}
The dimensions \(d_n\) grow with exponential rate equal to the
plastic constant, the positive root of \(x^3=x+1\).  Fourier
coefficients of a fixed K3 Borcherds product have their own
chamber-dependent asymptotics.  These sequences count objects in
distinct spaces.  A comparison of their growth is therefore attached
to a specified linear map from the Borcherds coefficient space to the
completed deformation complex.
\cite{Brown2011,GritsenkoNikulin1998}.
\end{remark}

\begin{remark}[The cyclic pro-obstruction]
\label{rem:nc-obs-infty}
For each cutoff \(m\), let \(\mathfrak c_{\leq m}\) be the cyclic
deformation complex truncated in motivic weight and let
\(\mathrm{obs}_{\leq m}\) be a represented obstruction cocycle.
A cyclic pro-obstruction consists of compatible transition maps
\[
 q_{m+1,m}\colon\mathfrak c_{\leq m+1}\longrightarrow
 \mathfrak c_{\leq m},
 \qquad
 q_{m+1,m}(\mathrm{obs}_{\leq m+1})
 =
 \mathrm{obs}_{\leq m}.
\]
When the \(q_{m+1,m}\) are degreewise surjective, the inverse system is
Mittag--Leffler and
\[
 \mathrm{obs}_\infty:=
 (\mathrm{obs}_{\leq m})_m
 \in\varprojlim_m\mathfrak c_{\leq m}
\]
is defined with \(\varprojlim^1=0\).  Brown's finite motivic
weight spaces and the finite-window topology provide the truncations;
the transition maps and their compatibility with cyclic closure are
the chain-level hypotheses.
\cite{Brown2011,MarklShniderStasheff2002}.
\end{remark}

\begin{remark}[Weight eleven: period coordinates and the order-eleven equation]
\label{rem:nc-phi11}
Brown's Hoffman basis theorem gives
\[
 d_{11}=d_9+d_8=5+4=9.
\]
The nine basis words are the four permutations of
\((2,3,3,3)\) and the five permutations of
\((2,2,2,2,3)\).  These words index the weight-$11$ motivic period
space.

Let \(\mathfrak c_{\mathrm{cyc}}\) be the completed cyclic deformation
dg Lie algebra and
\(\Theta_{<11}=\sum_{3\leq i<11}\hbar^i\psi_i\) a represented
Maurer--Cartan solution through order~$10$.  The order-$11$
obstruction is the closed cochain
\[
 R_{11}:=
 \frac12\!\sum_{\substack{i+j=11\\i,j\geq3}}[\psi_i,\psi_j].
\]
An order-$11$ term consists of a cyclic cochain
\(\psi_{11}\) satisfying
\[
 D\psi_{11}=-R_{11},
 \qquad
 \operatorname{Rot}_{11}(\psi_{11})=0.
\]
A chosen associator, a word-to-cochain map from the nine Hoffman
coordinates, and a comparison homotopy with the Borcherds scalar
channel supply the data entering this equation.  The dimension
\(d_{11}=9\) is the exact arithmetic input.
\cite{Brown2011,Hoffman1997,Drinfeld1990}.
\end{remark}

\begin{remark}[Weight twelve: the product-period sector]
\label{rem:nc-phi12}
Brown's dimension formula gives
\[
 d_{12}=d_{10}+d_9=7+5=12.
\]
Concretely, the Hoffman words consist of
\((3,3,3,3)\), the ten permutations of
\((2,2,2,3,3)\), and \((2,2,2,2,2,2)\).

The repeated period obeys the exact Newton identity
\begin{equation}
\label{eq:nc-zeta3333-newton}
 \zeta(3,3,3,3)
 =
 \frac{\zeta(3)^4-6\zeta(3)^2\zeta(6)+3\zeta(6)^2
 +8\zeta(3)\zeta(9)-6\zeta(12)}{24}.
\end{equation}
Indeed,
\[
 \sum_{r\geq0}\zeta(\underbrace{3,\ldots,3}_{r})t^r
 =
 \exp\!\left(
  \sum_{k\geq1}\frac{(-1)^{k-1}}{k}\zeta(3k)t^k
 \right),
\]
and extraction of \(t^4\) yields
\eqref{eq:nc-zeta3333-newton}.  Euler's formula places
\(\zeta^{\mathfrak m}(12)\) in
\(\mathbb Q\cdot(\zeta^{\mathfrak m}(2))^6\); hence every term on the
right belongs to the square of the positive-weight motivic ideal.
Thus
\[
 \operatorname{pr}_{\mathrm{ind}}
 \bigl(\zeta^{\mathfrak m}(3,3,3,3)\bigr)=0.
\]

A primitive weight-$12$ correction consists instead of a closed graph
cochain \(\gamma_{12}\), its action on
\(\mathfrak c_{\mathrm{cyc}}\), a period map, and homotopies comparing
the graph, KZ, and Borcherds images.  The Newton identity determines
the decomposable scalar sector of this comparison datum.
\cite{Brown2011,Brown2012,Hoffman1997}.
\end{remark}

\begin{remark}[Motivic quotients and cyclic inverse limits]
\label{rem:nc-ML-prolimit-12}
Let \(Z_{\leq m}^{\mathfrak m}\) denote the quotient of the motivic MZV
algebra by weights \(>m\).  The transition maps
\[
 Z_{\leq m+1}^{\mathfrak m}\twoheadrightarrow Z_{\leq m}^{\mathfrak m}
\]
are surjective by construction, and Brown's theorem gives finite
weight dimensions through the generating series
\((1-t^2-t^3)^{-1}\).  This inverse system is Mittag--Leffler.

A cyclic obstruction inverse system consists of completed deformation
complexes \(\mathfrak c_{\leq m}\), transition maps
\(q_{m+1,m}\), represented obstruction cocycles, and comparison maps
from \(Z_{\leq m}^{\mathfrak m}\).  Surjectivity of the
\(q_{m+1,m}\) implies the cyclic Mittag--Leffler condition and
\(\varprojlim^1\mathfrak c_{\leq m}=0\).  Thus the motivic quotient
system supplies one hypothesis; the cyclic transition maps supply the
remaining hypothesis for the pro-obstruction
\(\mathrm{obs}_\infty\).
\cite{Brown2011,MarklShniderStasheff2002}.
\end{remark}

\begin{remark}[Weight-twelve graph--associator comparison]
\label{rem:nc-grt1-weight-12}
A weight-$12$ primitive action on the completed deformation complex is
specified by the following chain data:
\[
 \gamma_{12}\in Z^0(\mathrm{GC}_2)_{12},\qquad
 \rho\colon\mathrm{GC}_2\longrightarrow
 \operatorname{Der}(\mathfrak c_{\mathrm{cyc}}),
\]
a period map on \([\gamma_{12}]\), and degree-\((-1)\) homotopies
intertwining \(\rho(\gamma_{12})\) with the chosen associator and
Borcherds comparison maps.  Under these hypotheses,
\(\rho(\gamma_{12})\) is the primitive weight-$12$ correction.
Willwacher's isomorphism
\(H^0(\mathrm{GC}_2)\cong\mathfrak{grt}_1\) identifies the ambient
Lie algebra of this action.  Equation
\eqref{eq:nc-zeta3333-newton} contributes zero to this primitive
coordinate and supplies its decomposable period sector.
\cite{Willwacher2015,Brown2012,Drinfeld1990}.
\end{remark}

\begin{remark}[The two weight-twelve automorphic domains]
\label{rem:nc-fake-monster-non-interference}
The K3 scalar section at order~$12$ is
\[
 \left(\frac{\Delta_5}{\eta^{12}}\right)^{12}
 =
 \frac{(\Phi_{10}^{\mathrm{un}})^6}{\eta^{144}},
\]
with raw modular weight \(-12\) and trivial
\(\Delta_5\)-character.  Its orthogonal domain is the K3 lattice of
signature \((2,1)\).  The Fake-Monster denominator \(\Phi_{12}\)
belongs to the orthogonal domain of
\(\mathrm{II}_{25,1}\), of signature \((25,1)\).
Accordingly the comparison complex records these scalar sections in
two typed summands.  The K3 order-$12$ comparison uses the first
summand.
\cite{Borcherds1998,GritsenkoNikulin1998}.
\end{remark}

\begin{remark}[Exact repeated-zeta scalar at weight twelve]
\label{rem:nc-c12-explicit}
Equation~\eqref{eq:nc-zeta3333-newton} gives
\[
 \zeta(3,3,3,3)
 =
 0.00029599901404317183550096081913354558\ldots
\]
and therefore
\[
 \frac{\zeta(3,3,3,3)}{12!}
 =
 6.1794994848278551783743690863150683\ldots\times10^{-13}.
\]
The first value follows from depth-one zeta values by Newton's
identity and independently from the increasing elementary-symmetric
sums
\[
 \sum_{N\geq n_1>n_2>n_3>n_4\geq1}
 \frac{1}{(n_1n_2n_3n_4)^3}.
\]
The factor \(12!\) is a scalar simplex normalisation.  Its occurrence
as a coefficient of a pentagon or cyclic cochain is determined by the
word-to-cochain map and the comparison homotopies of
Remark~\ref{rem:nc-grt1-weight-12}.
\cite{Hoffman1997,Brown2011}.
\end{remark}

\begin{remark}[Chain comparison at weight twelve]
\label{rem:nc-kz-depth-4-integral}
Let \(C_{\mathrm{KZ},12}\), \(C_{\mathrm{graph},12}\), and
\(C_{\mathrm{cyc},12}\) denote the weight-$12$ KZ-word,
graph-complex, and cyclic deformation complexes.  A chain comparison
is a diagram
\[
 C_{\mathrm{KZ},12}
 \xleftarrow{\ p_{12}\ }
 C_{\mathrm{graph},12}
 \xrightarrow{\ \rho_{12}\ }
 C_{\mathrm{cyc},12}
\]
together with a period map and a homotopy identifying the represented
KZ and cyclic images.  Once this datum is fixed,
\(\rho_{12}(\gamma_{12})\) defines the primitive cyclic class and
\(p_{12}(\gamma_{12})\) determines its chosen-associator coefficient.
The scalar of Remark~\ref{rem:nc-c12-explicit} occupies the
decomposable summand and has primitive projection zero.
\cite{Drinfeld1990,Willwacher2015,Brown2012}.
\end{remark}

\begin{remark}[Fourier and motivic comparison maps]
\label{rem:nc-fake-monster-fourier-check}
The K3 Borcherds product supplies a rational Fourier-coefficient
algebra, while the KZ associator supplies a motivic period algebra.
A map between these sources is part of the graph--trace comparison
datum.  In the direct-sum presentation of
Remark~\ref{rem:nc-fake-monster-non-interference}, the K3 scalar
section maps to the Borcherds trace summand and
\(\zeta(3,3,3,3)\) maps to the decomposable motivic summand.
A primitive cyclic term is supplied by
\(\rho_{12}(\gamma_{12})\).
\cite{Borcherds1998,GritsenkoNikulin1998,Brown2011}.
\end{remark}

\begin{remark}[Arithmetic and primitive data above weight twelve]
\label{rem:nc-open-horizon}
Brown's Hoffman basis theorem supplies the motivic weight dimensions
in every weight through
\[
 \sum_{n\geq0}d_nt^n=\frac{1}{1-t^2-t^3}.
\]
Depth-graded primitive graph classes form a second problem.  At each
weight \(n\), a primitive correction is specified by a represented
class \(\gamma_n\in Z^0(\mathrm{GC}_2)_n\), its period, its action on
the deformation complex, and the comparison homotopies.  The
Broadhurst--Kreimer series organises conjectural ranks for these
depth-graded pieces, while the Hoffman basis controls the proved
period-algebra dimensions.
\cite{Brown2011,Brown2012,BroadhurstKreimer1997}.
\end{remark}

\begin{remark}[Typed asymptotic comparison]
\label{rem:nc-asymptotic-11-12}
The motivic sequence \(d_n\) counts Hoffman period coordinates.  The
Fourier coefficients of the K3 Borcherds product count root
multiplicities after a chamber and a Fourier direction have been
fixed.  These sequences inhabit distinct vector spaces.  An
asymptotic comparison is therefore a theorem about a specified linear
map between the period and root-multiplicity spaces.  The present
chapter retains the exact values
\[
 d_{11}=9,\qquad d_{12}=12,
\]
and places the Borcherds coefficients in their scalar comparison
summand.
\cite{Brown2011,GritsenkoNikulin1998}.
\end{remark}

\begin{remark}[Target transitivity for the K$3$-BKM comparison datum]
\label{rem:nc-grt1-bkm-transitivity-vol1}
The Drinfeld twist action of the pro-unipotent
Grothendieck--Teichm\"uller group
$\exp(\widehat{\mathfrak{grt}}_1)$ on Drinfeld associators is
transitive on the affine Kac--Moody sub-locus by
Etingof--Kazhdan~$2000$ Part~V Thm.~$2.2$. For the K$3$-BKM target
datum of $\mathfrak g_{\Delta_5}$ the imaginary cone contributes an
extra class
\[
\mathrm{ob}^{\mathrm{GN}} \in
H^2(\mathfrak{grt}_1;\widehat{\mathrm{Imag}}_{\Delta_5}) .
\]
Vol.~I controls the affine EK mechanism and the motivic associator
coordinates. It does not prove the vanishing of
$\mathrm{ob}^{\mathrm{GN}}$, and Brown's unconditional motivic basis
through weight~$12$ is not a classification of the compact K$3$-BKM
Hall source through weight~$12$. Under the target-side hypothesis
$\mathrm{ob}^{\mathrm{GN}}=0$ on the Koszul locus
$\overline{\mathcal A_2}\smallsetminus\bigcup_{n\;\text{admissible}}H_n$,
the super-Etingof--Kazhdan target quantisation space
$\Quant^{\mathrm{GN, Koszul}}_{\mathrm{tgt}}(\mathfrak g_{\Delta_5})/
(\mathbb Z/2)_{\mathrm{super}}$ becomes a torsor over
$\exp(\widehat{\mathfrak{grt}}_1)$. Recognition of a compact Hall
source $\mathbf H_{\Delta_5}$ requires the Vol.~III source gates
together with the cohomology/Heegner-characteristic comparison.
\cite{Drinfeld1990,Brown2012,Willwacher2015,SchnepperWillwacher2017,
 GritsenkoNikulin1998,Borcherds1995}.
\end{remark}

\begin{theorem}[Conditional target $\mathrm{GRT}_1^{\mathrm{super}}$-torsor
for $\mathfrak g_{\Delta_5}$]
\label{thm:nc-grt1-super-torsor-delta5}
\ClaimStatusConditional
Let
\[
\mathcal Q^{\mathrm{tgt}}(\mathfrak g_{\Delta_5})
\;:=\;
\bigl\{\mathbf H : \mathbf H/\hbar = U(\mathfrak g_{\Delta_5}),
\,\mathbf H \text{ super-quasi-Hopf},
\,\mathrm{GN}\text{-compatible}\bigr\}\big/
\text{super-quasi-Hopf gauge}
\]
denote the space of super-Etingof--Kazhdan quantisations of the
K$3$-BKM Lie superalgebra $\mathfrak g_{\Delta_5}$ compatible with
the Gritsenko--Nikulin automorphic cobracket $\delta_{\mathrm{GN}}$,
modulo super-quasi-Hopf gauge equivalence, restricted to the Koszul
locus
$\overline{\mathcal A_2}\smallsetminus\bigcup_{n\;\text{admissible}}H_n$
of Theorem~B. Let
\[
\mathrm{GRT}_1^{\mathrm{super}}
\;:=\;
\exp\bigl(\widehat{\mathfrak{grt}}_1\bigr)\rtimes(\mathbb Z/2)_{\mathrm{super}}
\]
denote the semidirect product of the pro-unipotent Grothendieck--Teichm\"uller
group with the $\mathbb Z/2$ super-sign of the Etingof--Kazhdan
super-twist. Assume:
\begin{enumerate}
\item the Etingof--Kazhdan super-quantisation functor extends from the
affine Kac--Moody Manin-pair backbone to the completed
Gritsenko--Nikulin BKM target datum;
\item the imaginary-cone obstruction
$\mathrm{ob}^{\mathrm{GN}}\in
H^2(\mathfrak{grt}_1;\widehat{\mathrm{Imag}}_{\Delta_5})$ vanishes on
the Koszul locus;
\item whenever $\mathcal Q^{\mathrm{tgt}}(\mathfrak g_{\Delta_5})$ is
used to recognise a compact object $\mathbf H_{\Delta_5}$, the
Vol.~III Hall/CoHA source-recognition gates and the
cohomology/Heegner-characteristic comparison have been supplied.
\end{enumerate}
Under these hypotheses
$\mathcal Q^{\mathrm{tgt}}(\mathfrak g_{\Delta_5})$ is a principal
homogeneous space over $\mathrm{GRT}_1^{\mathrm{super}}$: any two
target super-EK quantisations are related by a unique
$\mathrm{GRT}_1^{\mathrm{super}}$ Drinfeld twist. Brown's motivic
basis gives unconditional associator coordinates through weight~$12$;
it does not give an unconditional classification of the K$3$-BKM
compact Hall source through weight~$12$.
\end{theorem}

\begin{proof}
Three ingredients assemble into the conditional target statement.

\emph{Ingredient~$1$ (affine backbone, Etingof--Kazhdan~$1996$--$2000$
Parts~I--V).} By EK~$1996$ Part~I Thm.~$4.2$ the pro-unipotent group
$\exp(\widehat{\mathfrak{grt}}_1)$ acts freely and transitively on
the space of Drinfeld associators $\Phi_{\mathrm{KZ}}$ associated to
a Manin pair; the super-extension (EK~$2000$ Part~V Thm.~$2.2$) adds
the $\mathbb Z/2$-graded module structure. On the affine Kac--Moody
sub-locus $\Delta_{\mathrm{aff}}\subset\Delta(\mathfrak g_{\Delta_5})$
(real-root part, signature $(2, 1)$) this
recovers the classical affine-KM torsor statement. This is the
unconditional Vol.~I mechanism.

\emph{Ingredient~$2$ (conditional imaginary-cone extension).} On
$\Delta_+^{\mathrm{im}}$ the multiplicities $\mathrm{mult}(\alpha) =
c_{K3}(4nm - \ell^2)$ obey the pro-completion convergence of
Definition~$\ref{V3-def:qgf-imag-pro}$; the $\grt_1$-action lifts to
$\widehat{\Imag}$ through the Ihara-bracket filtration
(Drinfeld~$1990$~\S$5$; Furusho~$2011$ \emph{Ann.\ Math.}~$174$
Thm.~$4$). The obstruction cocycle
$\mathrm{ob}^{\mathrm{GN}} \in H^2(\grt_1;\widehat\Imag)$ of
Remark~\ref{rem:nc-grt1-bkm-transitivity-vol1} is the additional
target hypothesis. If it vanishes, the affine torsor transports to the
completed Gritsenko--Nikulin target datum; without that vanishing,
Vol.~I has no transitivity theorem for the K$3$-BKM target.

\emph{Ingredient~$3$ (freeness).} Equivariance of the
Gritsenko--Nikulin character identity
(Definition~$\ref{V3-def:qgf-quant-delta5}$(c)) under
$\exp(\widehat{\grt}_1)$-twist is the content of
Theorem~$\ref{V3-thm:qgf-grt1-bkm-transitivity}$(i). The fixed-point set
is trivial because the Ihara bracket is free on the generators
$\sigma_3, \sigma_5, \sigma_7, \ldots$ (Brown~$2012$ Thm.~$1.2$;
graph-complex freeness is Willwacher~$2015$ Thm.~$1.1$).
Non-triviality: Drinfeld~$1990$ \S$5$ Remark~$3$ shows that the
conjugate associator $f\cdot\Phi_{\mathrm{KZ}}\cdot f^{-1}$ differs
from $\Phi_{\mathrm{KZ}}$ when $f$ is any non-identity element of
$\exp(\widehat\grt_1)$. The super-sign factor of
$(\mathbb Z/2)_{\mathrm{super}}$ acts non-trivially on the braiding
by EK~$2000$ Part~V Prop.~$2.1$; together the semidirect product
acts freely.

Combining the three ingredients, on the Koszul locus the
$\mathrm{GRT}_1^{\mathrm{super}}$-action on
$\mathcal Q^{\mathrm{tgt}}(\mathfrak g_{\Delta_5})$ is simply
transitive precisely under the stated target obstruction hypothesis.
The weight~$12$ assertion belongs only to the motivic-coordinate
window: Brown~$2011$ \emph{Ann.\ Math.}~$175$ Thm.~$1.2$ supplies the
MZV basis through weight~$12$; above weight~$12$ those coordinates are
conditional on Zagier--Hoffman. Neither fact constructs nor classifies
the compact Hall source $\mathbf H_{\Delta_5}$.
\end{proof}

\begin{remark}[Associator orbit and weight filtration]
\label{rem:nc-grt1-orbit-description}
Fix a reference Drinfeld associator
\(\Phi^{\mathrm{KZ}}_{\mathrm{ref}}\).  Its orbit is
\[
 \mathrm{GRT}_1\cdot\Phi^{\mathrm{KZ}}_{\mathrm{ref}}
 =
 \{\,g\cdot\Phi^{\mathrm{KZ}}_{\mathrm{ref}}:
      g\in\mathrm{GRT}_1\,\}.
\]
The weight filtration on \(\mathfrak{grt}_1\) induces the corresponding
filtration on this orbit.  At weight~$12$, a primitive coordinate is
defined by a represented graph class \(\gamma_{12}\) and the
graph--associator comparison of
Remark~\ref{rem:nc-grt1-weight-12}.  The period
\(\zeta(3,3,3,3)\) belongs to the decomposable summand by
Equation~\eqref{eq:nc-zeta3333-newton}.  The action on the K3
Borcherds trace is supplied by the continuous imaginary-cone
representation and its intertwining homotopy.
\cite{Drinfeld1990,Willwacher2015,Brown2012}.
\end{remark}

\begin{remark}[Target Maurer--Cartan moduli]
\label{rem:nc-grt1-geometric-interpretation}
The completed super-convolution dg Lie algebra
\[
 \mathfrak c_{\Delta_5}:=
 \mathrm{Conv}\bigl(
 B^{\mathrm{ch}}(\mathfrak g_{\Delta_5}^{\mathrm{super}}),
 \Omega^{\mathrm{ch}}\bigr)
\]
defines a formal Maurer--Cartan moduli problem.  A continuous
\(\mathfrak{grt}_1\)-action on \(\mathfrak c_{\Delta_5}\), together
with the imaginary-cone cochain comparison, induces an action on this
formal moduli problem.  For a represented graph class
\(\gamma_n\), the tangent vector is
\(\rho(\gamma_n)\).  Its associator coefficient is obtained through
the chosen period map.  At weight~$12$, the repeated-zeta scalar has
zero primitive projection, while the primitive tangent vector is
\(\rho(\gamma_{12})\).
\cite{Kontsevich2003,Willwacher2015,Drinfeld1990}.
\end{remark}

\begin{corollary}[$\Delta_5$ target criterion for $\mathbf H_{\Delta_5}$]
\label{cor:nc-h-delta5-unique-up-to-grt1}
\ClaimStatusConditional
Assume the Vol.~III source-recognition package for
$\mathbf H_{\Delta_5}$: finite Hall/CoHA source, compact Hopf pairing,
source-to-target comparison maps, radical descent, no-extra-relations,
PBW compatibility, and exact inverse-limit passage
(Vol.~III Chapter~\ref{V3-ch:k3-chiral-bialgebra-platonic},
Definition~\ref{V3-def:k3-chiral-bialgebra} and
Theorem~\ref{V3-thm:k3-structural-genesis}). Assume also a proved
cohomology/Heegner-characteristic comparison
\[
\chi_{\mathrm{Heg}}\colon
H^2(\mathfrak g_{\Delta_5})^{\mathbb Z/2,\,K(1)}
\longrightarrow
\mathbb C\cdot\Delta_5
\]
identifying the recognised source deformation class with the
Bruinier--Heegner characteristic line and proving that this target line
is the relevant one-dimensional comparison quotient. Under these
hypotheses the $\Delta_5$-line is a complete target comparator for the
$\mathrm{GRT}_1^{\mathrm{super}}$-torsor of
Theorem~\ref{thm:nc-grt1-super-torsor-delta5}:
\[
\mathbf H_{\Delta_5}^{\mathrm{tgt}}
\;\in\;
\mathcal Q^{\mathrm{tgt}}(\mathfrak g_{\Delta_5})
\big/\mathrm{GRT}_1^{\mathrm{super}}
\quad\text{with}\quad
\mathrm{im}(\chi_{\mathrm{Heg}})=\mathbb C\cdot\Delta_5 .
\]
Equivalently, after those gates are supplied, any two recognised
compact sources whose Heegner characteristic has the same normalised
$\Delta_5$ image lie in the same
$\mathrm{GRT}_1^{\mathrm{super}}$-orbit. Without the source-recognition
and Heegner-comparison inputs, Vol.~I retains only this target criterion.
\end{corollary}

\begin{proof}
Theorem~\ref{thm:nc-grt1-super-torsor-delta5} controls the target
super-EK quantisations of the BKM Manin-pair datum; it neither
constructs the compact Hall source nor computes
$H^2(\mathfrak g_{\Delta_5})$ intrinsically. The Vol.~III recognition
package promotes the target datum to a sourced compact object only
after the Hall/CoHA source, pairing, comparison, PBW, and inverse-limit
gates have been verified. The additional
cohomology/Heegner-characteristic comparison identifies the source
deformation class with the automorphic $\Delta_5$ line. Under precisely
these assumptions, equality of the normalised $\Delta_5$ characteristic
places two recognised sources in the same target
$\mathrm{GRT}_1^{\mathrm{super}}$-orbit by the torsor theorem. This is a
conditional recognition criterion, not an unconditional classification
of $\mathbf H_{\Delta_5}$ inside Vol.~I.
\end{proof}

\begin{remark}[Assembly of the conditional target torsor proof from EK Parts I--V]
\label{rem:nc-grt1-ek-parts-assembly}
The proof of Theorem~\ref{thm:nc-grt1-super-torsor-delta5} assembles
five Etingof--Kazhdan structural pieces into a conditional target
transitivity statement. Each piece is classical for quasi-triangular
Lie bialgebras; the K$3$-BKM passage enters only through the completed
Gritsenko--Nikulin target datum and the imaginary-cone vanishing
hypothesis of Theorem~\ref{thm:nc-grt1-super-torsor-delta5}. The
decomposition makes the target mechanism verifiable against five
independent primary witnesses:
\begin{itemize}
\item \emph{Part~I (EK~$1996$ Sel.~Math.~$2$).} Existence and
 uniqueness, up to $\grt_1$-gauge, of a quantisation
 $U_\hbar(\mathfrak a)$ for any quasi-triangular Lie bialgebra
 $(\mathfrak a, \delta)$. Applied here to the affine sub-lattice
 $\Delta_{\mathrm{aff}} \subset \Delta(\mathfrak g_{\Delta_5})$
 (signature $(2, 1)$, real-root part), supplying the affine
 Kac--Moody backbone with $\delta$ equal to the restriction of
 $\delta_{\mathrm{GN}}$; the torsor structure on the affine
 quantisation moduli is Corollary~\ref{V3-cor:qgf-grt1-affine-recovery}.
\item \emph{Part~II (EK~$1998$ Sel.~Math.~$4$~\S$2$).} Functoriality
 of the EK functor with respect to Manin-pair morphisms. Applied here
 to the Manin-pair embedding
 $(\Delta_{\mathrm{aff}}, \mathrm{triangular})
 \hookrightarrow (\Delta(\mathfrak g_{\Delta_5}), \mathrm{imaginary})$;
 under the obstruction-vanishing hypothesis, this functoriality
 transports the affine torsor structure to the target imaginary-cone
 extension along the Gritsenko--Nikulin lift.
\item \emph{Part~III (EK~$1998$ Sel.~Math.~$4$~\S$3$).} Twist
 invariance: if $\Phi, \Phi'$ are two Drinfeld associators related
 by $f \in \exp(\widehat\grt_1)$, the induced quantisations
 $Q_\Phi, Q_{\Phi'}$ are canonically isomorphic as super-quasi-Hopf
 algebras via the twist $R^{\Phi'} = f_{321} R^\Phi f^{-1}$. This
 supplies the \emph{action} of $\mathrm{GRT}_1$ on
 $\mathcal Q^{\mathrm{tgt}}(\mathfrak g_{\Delta_5})$ by Drinfeld
 twist; together with Part~I and the target obstruction hypothesis it
 supplies transitivity.
\item \emph{Part~IV (EK~$2000$ Sel.~Math.~$6$).} Kac--Moody
 extension: the affine-KM case is controlled by the imaginary
 sector being trivial. The $\mathfrak g_{\Delta_5}$-case is its
 hyperbolic extension; the imaginary-cone multiplicities
 $\mathrm{mult}(\alpha) = c_{K3}(4nm - \ell^2)$ (EOT~$2011$) replace
 the zero multiplicities of the affine case.
\item \emph{Part~V (EK~$2000$ Sel.~Math.~$6$~\S$2$).} Super-extension:
 $\mathbb Z/2$-graded Lie bialgebras admit super-EK quantisations,
 with the super-sign entering as an object-level grading that
 commutes with the associator-level $\grt_1$-twist of Part~III.
 This supplies the $(\mathbb Z/2)_{\mathrm{super}}$ semidirect factor
 of $\mathrm{GRT}_1^{\mathrm{super}}$ and upgrades the Part~III twist
 to a super-twist. On the K$3$-BKM side, the super-sign is the
 fermion-number grading of the $\mathcal N = 4$ elliptic genus
 (EOT~$2011$; Cheng--Duncan--Harvey~$2014$ umbral moonshine).
\end{itemize}
Ingredients~$1$--$3$ of the proof of
Theorem~\ref{thm:nc-grt1-super-torsor-delta5} correspond to
Parts~(I, IV)+Part~V (affine + BKM + super), Parts~(II, III)+IV
(functoriality + twist + imaginary cone), and Part~(III)+V
(free action + super-sign non-triviality), respectively.
\cite{EtingofKazhdan1996,EtingofKazhdan1998a,EtingofKazhdan1998b,
 EtingofKazhdan2000,EtingofKazhdan2008}.
\end{remark}

\begin{remark}[Target verification via $\phi^{(3)}$ transforming under
$\sigma_3 \in \widehat{\mathfrak{grt}}_1$]
\label{rem:nc-grt1-phi3-verification}
A concrete chain-level witness of the conditional target
$\mathrm{GRT}_1^{\mathrm{super}}$ action is the weight-$3$ component
$\phi^{(3)}$ of the
universal associator, written in the reference basis as
\[
\phi^{(3)}
\;=\;
\zeta(3)\,c_{\mathrm{symm}}^{(3)}
\;+\; c_{\mathrm{KM}}^{(3)}\,\zeta(3)\cdot\Omega_{\mathrm{KM}}
\;+\; c_{\Phi_{10}^{\mathrm{un}}/\eta^{24}}^{(3)}\cdot[\Phi_{10}^{\mathrm{un}}/\eta^{24}]_3
\;+\; c_{\mathrm{skew}}^{(3)}\cdot\zeta(3)\,t_{\mathrm{skew}}
\]
(cf.~Vol~III Rem.~\ref{rem:qgf-phi3-explicit}), where
$c_{\mathrm{symm}}^{(3)}, c_{\mathrm{KM}}^{(3)},
c_{\Phi_{10}^{\mathrm{un}}/\eta^{24}}^{(3)}, c_{\mathrm{skew}}^{(3)} \in \mathbb Q$
are the four reference coefficients. The generator
$\sigma_3 \in \widehat{\mathfrak{grt}}_1$ (Ihara weight~$3$) acts by
\[
\sigma_3\cdot\phi^{(3)}
\;=\;
\phi^{(3)} + \zeta(3)\cdot\iota_3
\]
with $\iota_3 \in \mathfrak g_{\Delta_5}^{\otimes 3}$ the Ihara
infinitesimal tangent vector at weight~$3$, which in the present
basis shifts only the $c_{\mathrm{symm}}^{(3)}$ coefficient (the
depth-$1$ MZV class), leaves the Kac--Moody coefficient
$c_{\mathrm{KM}}^{(3)}\,\zeta(3)\,\Omega_{\mathrm{KM}}$ gauge-equivalent
through a conjugation by $\exp(\hbar\cdot\Omega_{\mathrm{KM}})$
(the free-field cocycle of the Heisenberg subalgebra), and fixes
the Borcherds leg $[\Phi_{10}^{\mathrm{un}}/\eta^{24}]_3$ pointwise (the imaginary
cone decouples from weight~$3$: the smallest non-trivial imaginary
contribution to $[\Phi_{10}^{\mathrm{un}}/\eta^{24}]_n$ is at $n = 10$ by the
Gritsenko weight, and the weight-$3$ slice of the imaginary-cone
bivector vanishes by Hardy--Ramanujan asymptotics truncated at
$N \le 3$).

At weight~$5$, the generator $\sigma_5$ produces a depth-$1$ shift
$\sigma_5\cdot\phi^{(5)} = \phi^{(5)} + \zeta(5)\cdot\iota_5$, again
fixing the Borcherds leg; the first non-trivial interaction with
the imaginary cone is at weight~$10$ via the Ihara--Zagier relation
$\zeta(10) = \frac{4}{17}[\zeta(3)\zeta(7) + \zeta(5)^2]$
(which is fixed by $\grt_1$; the depth-$1$ odd sector at weight~$10$
is zero).  At
weight~$12$, Equation~\eqref{eq:nc-zeta3333-newton} supplies the
decomposable repeated-zeta sector and gives zero after primitive
projection.  A primitive correction is the image
$\rho(\gamma_{12})$ of a represented graph class under the
comparison datum of Remark~\ref{rem:nc-grt1-weight-12}.  The affine
limit, the period map, and the graph-complex action are the three
inputs that determine the corresponding weight-$12$ twist.
\cite{Drinfeld1990,Brown2012,EtingofKazhdan2000,GritsenkoNikulin1998}.
\end{remark}

\begin{remark}[Filtered associator orbit]
\label{rem:nc-grt1-explicit-orbit}
For a reference associator
\(\Phi^{\mathrm{KZ}}_{\mathrm{ref}}\), the orbit map is
\[
 \mathrm{GRT}_1\longrightarrow\operatorname{Assoc},
 \qquad
 g\longmapsto g\cdot\Phi^{\mathrm{KZ}}_{\mathrm{ref}}.
\]
The weight filtration on \(\mathfrak{grt}_1\) filters its tangent
space.  Drinfeld's odd classes give distinguished tangent directions
in odd weights.  A weight-$12$ primitive direction is represented by
a graph cocycle \(\gamma_{12}\) together with the comparison datum of
Remark~\ref{rem:nc-grt1-weight-12}.  Its cyclic action is
\(\rho(\gamma_{12})\).  The repeated-zeta value of
Remark~\ref{rem:nc-c12-explicit} has zero primitive projection and
therefore belongs to the product-period coordinate ring of the orbit.
\cite{Drinfeld1990,Willwacher2015,Brown2012}.
\end{remark}

\begin{remark}[Motivic period map for the target action]
\label{rem:nc-grt1-motivic-geometric}
The unipotent motivic fundamental group of
\(\mathbb P^1\smallsetminus\{0,1,\infty\}\) supplies the period
realisation of the chosen associator.  Willwacher's isomorphism
\[
 H^0(\mathrm{GC}_2)\cong\mathfrak{grt}_1
\]
supplies the graph-complex tangent Lie algebra.  A target action on
the K3-BKM deformation complex is consequently determined by a
continuous representation
\[
 \rho_{\Imag}\colon\mathfrak{grt}_1
 \longrightarrow
 \operatorname{Der}_{\mathrm{cont}}(\widehat{\Imag}),
\]
a graph-to-Chevalley--Eilenberg cochain map, and homotopies
intertwining the period, cyclic, and Borcherds images.  In weight~$12$
the primitive action is the image of the represented class
\(\gamma_{12}\); Equation~\eqref{eq:nc-zeta3333-newton} determines
the decomposable period component.
\cite{DeligneGoncharov2005,Drinfeld1990,Willwacher2015,Brown2012}.
\end{remark}

\section{Deligne--Tamarkin Massey triple $\langle r, r, r\rangle$ at $E_8$}
\label{sec:nc-massey-triple-rrr}

\begin{proposition}[\(E_8\) cubic invariant channels]
\label{prop:nc-massey-triple-rrr-E8}
\ClaimStatusProvedHere
Let \(\mathfrak g\) be a simple Lie algebra over a field of
characteristic zero with invariant form \(\langle-,-\rangle\).
For any invariant cubic tensor
\(\tau\in(\mathfrak g^{\otimes3})^{\mathfrak g}\), its symmetric and
alternating projections satisfy
\[
 \pi_{\mathrm{sym}}\tau\in
 \operatorname{Sym}^3(\mathfrak g)^{\mathfrak g},
 \qquad
 \pi_{\mathrm{alt}}\tau\in
 (\Lambda^3\mathfrak g)^{\mathfrak g}.
\]
For \(\mathfrak g=E_8\),
\[
 \operatorname{Sym}^3(E_8)^{E_8}=0,
 \qquad
 (\Lambda^3E_8)^{E_8}
 =\bk\cdot\bigl[(x,y,z)\mapsto\langle x,[y,z]\rangle\bigr].
\]
Hence the symmetric projection of every invariant cubic local tensor
for \(V_k(E_8)\) is zero, while its alternating projection is a scalar
multiple of the Cartan three-cocycle.
\end{proposition}

\begin{proof}
Chevalley's theorem identifies
\(\operatorname{Sym}(\mathfrak g)^{\mathfrak g}\) with a polynomial
algebra whose generator degrees are the exponents plus one.  For
\(E_8\) these degrees are
\[
 2,8,12,14,18,20,24,30.
\]
Thus the degree-three invariant space is zero.  For every simple
\(\mathfrak g\), the invariant alternating three-forms form the
one-dimensional primitive degree-three Lie algebra cohomology class,
represented by the Cartan cocycle
\((x,y,z)\mapsto\langle x,[y,z]\rangle\).
\end{proof}

\begin{remark}[$E_8$ invariant-theoretic splitting]
\label{rem:nc-E8-vanishing-cyclic-associator}
Chevalley's degree list gives
\(\operatorname{Sym}^3(E_8)^{E_8}=0\).  The alternating invariant
belongs to a different summand:
\[
 (\Lambda^3E_8)^{E_8}
 =\bk\cdot\bigl[(x,y,z)\longmapsto\langle x,[y,z]\rangle\bigr].
\]
Thus the symmetric projection of the cubic local tensor is zero,
while the alternating projection is the Cartan three-cocycle up to
the normalization of the invariant form.  A triple Massey package
consists of a cocycle representative, explicit null-homotopies for its
pairwise products, and the quotient by the resulting indeterminacy.
The classical Yang--Baxter equation supplies bracket compatibility;
the null-homotopies are additional chain data.  A scalar residue then
requires \(\mathsf H_{\mathrm{res}}\), while a KZ period such as
\(\zeta(3)\) requires the word-to-cochain comparison in
\(\mathsf H_{\mathrm{mot}}\).  The Sugawara conformal vector compares
the local Virasoro subalgebra, and a comparison of cubic residue
scalars is an additional chain map.
\cite{Kohno1985,Sugawara1968,BarNatan1995}.
\end{remark}

\section{Motivic associator coefficients above weight \(12\)}
\label{sec:nc-brown-above-12-double-conditional}

The motivic multiple-zeta algebra and the pentagon deformation complex
are distinct objects.  A coefficient \(\phi^{(n)}\) is obtained from a
motivic period only after a word-to-cochain comparison has been
constructed.  We denote this package by \(\mathsf H_{\mathrm{mot}}\).
It consists of a filtered chain map from the motivic associator
complex, representatives for the chosen depth-graded classes, and
homotopies intertwining the pentagon differential.

\begin{proposition}[Motivic comparison problem for \(\phi^{(n)}\)]
\label{prop:nc-double-conditional-phi-n}
\ClaimStatusOpen
\emph{Type signature:}
\textup{(Open quadrant, completed pentagon deformation presentation,
Beilinson level~\(2\), \(\mathsf H_{\mathrm{mot}}\)).}
For \(n\ge13\), the construction of \(\phi^{(n)}\) from motivic
multiple zeta values requires:
\begin{enumerate}[label=\textup{(M\arabic*)}]
\item a represented weight-\(n\) cocycle in the motivic associator
 complex;
\item a filtered word-to-pentagon cochain map;
\item a comparison homotopy identifying its coboundary with the
 arity-\(n\) deformation obstruction;
\item a choice of basis and integral normalization for every
 depth-graded quotient used in the scalar projection.
\end{enumerate}
The Zagier--Hoffman and Broadhurst--Kreimer conjectures may supply
dimension and spanning information for the source.  The package
\(\mathsf H_{\mathrm{mot}}\) supplies the chain-level comparison to
\(\phi^{(n)}\).
\end{proposition}

\begin{remark}[The weight-\(25\) test]
\label{rem:nc-phi-25-tightening}
Weight \(25\), depth \(5\), is a concrete test of the complete
package: one must choose a represented motivic class, compute its
pentagon image, and compare it with the lower-depth filtration.  A
Borcherds automorphic form belongs to a further target lane and enters
after a separate automorphic comparison morphism has been specified.
\end{remark}

\begin{proposition}[The depth-five iterated integral at weight \(25\)]
\label{prop:nc-phi-25-depth-5-coefficient}
\ClaimStatusConditional
The convergent iterated integral represented by the word
\[
 \omega_1\omega_0^2\,
 \omega_1\omega_0^2\,
 \omega_1\omega_0^4\,
 \omega_1\omega_0^4\,
 \omega_1\omega_0^8
\]
is the numerical multiple zeta value
\[
 \zeta(3,3,5,5,9).
\]
Assume the Broadhurst--Kreimer rank-one prediction for
\(\operatorname{gr}^{\mathcal D}_5\mathcal L_{\mathrm{mot}}^{(25)}\),
a nonzero projection of this class to that quotient, and
\(\mathsf H_{\mathrm{mot}}\) with simplex normalization
\(\Gamma(26)^{-1}\).
Then the corresponding scalar coordinate of the represented pentagon
cochain is
\[
 \frac{\zeta(3,3,5,5,9)}{\Gamma(26)}
 \quad\text{modulo the coordinates of the depth-\(\le4\) filtration.}
\]
\end{proposition}

\begin{proof}
The iterated-integral identity is the standard word realization of a
convergent multiple zeta value.  The rank-one hypothesis identifies its
nonzero depth-five projection up to scalar.  The represented
word-to-cochain map and the specified simplex normalization fix that
scalar, while the quotient records the lower-depth ambiguity.
\end{proof}

\begin{remark}[First complete depth-five comparison]
\label{rem:nc-phi-25-first-double-conditional}
The proposition separates three verifications: the numerical period,
its depth-graded motivic class, and its image in the pentagon complex.
The first is classical.  The second is governed by the stated
depth-graded hypotheses.  The third is the chain-level content of
\(\mathsf H_{\mathrm{mot}}\).
\end{remark}

\section{Conditional GRT$_1$-transitivity on full BKM target loci}
\label{sec:nc-grt1-full-bkm-conditional}

\begin{theorem}[Conditional GRT$_1^{\mathrm{super}}$-transitivity on full BKM target loci]
\label{thm:nc-grt1-full-bkm-conditional}
\ClaimStatusConditional
Let $\mathcal Q_{\mathrm{full}}(\mathfrak g_{\mathrm{BKM}})$ be the
space of super-EK quantisations of an arbitrary \textup{(}non-affine,
non-K$3$-specialised\textup{)} Borcherds--Kac--Moody Lie superalgebra
$\mathfrak g_{\mathrm{BKM}}$ with Cartan of rank $r$ and signature
$(p, q)$, $p + q = r$, modulo super-quasi-Hopf gauge. The
$\mathrm{GRT}_1^{\mathrm{super}}$-action by Drinfeld twist on
$\mathcal Q_{\mathrm{full}}$ is transitive, conditionally on:
\textup{(a)~vanishing of the BKM imaginary-cone obstruction
$H^2(\mathfrak{grt}_1; \widehat{\mathrm{Imag}}_{\mathrm{BKM}}) = 0$
on the full Koszul locus of} $\mathfrak g_{\mathrm{BKM}}$; and
\textup{(b)~Zagier--Hoffman + Broadhurst--Kreimer at every weight
$\le$ the motivic depth of the Cartan type}. The affine case
\textup{(}$\mathfrak g_{\mathrm{BKM}} = \widehat{\mathfrak g}^{\mathrm{aff}}$,
imaginary cone trivial\textup{)} is unconditional
\textup{(}Etingof--Kazhdan~$2000$ Part~V Thm.~$2.2$\textup{)}. The
K$3$-specialised target datum $\mathfrak g_{\Delta_5}$ is the
conditional case of
Theorem~\ref{thm:nc-grt1-super-torsor-delta5}: its motivic associator
coordinates are unconditional through weight~$12$, but the compact
Hall-source recognition and any classification of
$\mathbf H_{\Delta_5}$ require the Vol.~III gates plus the
cohomology/Heegner comparison. The general BKM case remains
conjectural pending the imaginary-cone obstruction computation, which
is an open question of Gritsenko--Nikulin type.
\end{theorem}

\begin{proof}[Proof of the conditionality]
Three independent obstructions separate the full-BKM case from the
K$3$-specialised one: (i)~higher imaginary-cone multiplicities
\textup{(}Borcherds $1992$ \emph{Invent.\ Math.}~$109$ Thm.~$10.1$
gives only the existence of the denominator; explicit multiplicities
for a generic Lorentzian lattice of rank~$> 3$ are unknown\textup{)};
(ii)~absence of a universal modular covering analogous to the
Siegel paramodular $K(1)$ \textup{(}each rank-$r$ Lorentzian lattice
admits its own automorphic cover, and the GRT$_1$-action does not
descend uniformly\textup{)}; (iii)~the obstruction cocycle
$\mathrm{ob}^{\mathrm{BKM}} \in H^2(\mathfrak{grt}_1; \widehat{\mathrm{Imag}}_{\mathrm{BKM}})$
need not vanish generically: on $\mathrm{II}_{25, 1}$ the Fake Monster
BKM $\mathfrak m_{\mathrm{FakeMonster}}$ carries a distinct vanishing
mechanism \textup{(}Borcherds~$1998$ Thm.~$13.1$\textup{)} from the
K$3$-BKM, and the two are not related by a universal argument. The
three obstructions specialise: (i) is trivial in the affine case
(no imaginary roots); (ii) is satisfied for K$3$-BKM via the Igusa
$\Phi_{10}^{\mathrm{un}}$ cover; (iii) becomes the K$3$ target-side
vanishing hypothesis of
Theorem~\ref{thm:nc-grt1-super-torsor-delta5}. Vol.~I therefore proves
the affine EK and motivic mechanisms, not the source classification of
$\mathbf H_{\Delta_5}$. A general proof requires a uniform
imaginary-cone vanishing argument across all Lorentzian Kac--Moody root
systems, which is an open problem of Gritsenko--Nikulin type
\textup{(}Gritsenko--Nikulin~$2002$ \emph{Russian Math.\ Surveys}~$57$
\S$5$ open question~$7$\textup{)}.
\cite{Borcherds1992,Borcherds1998,GritsenkoNikulin1998,
 GritsenkoNikulin2002,EtingofKazhdan2000}.
\end{proof}

\begin{remark}[Landscape-row scope]
\label{rem:nc-grt1-status-per-row}
The conditionality of Theorem~\ref{thm:nc-grt1-full-bkm-conditional}
reads per archetype.  Rows $\bG/\bL/\bCc$ lie in the affine EK lane.
Row $\bM$ uses the motivic residue package
$\mathsf H_{\mathrm{mot}}$ and its Tate-factorization package
$\mathsf H_{\mathrm{Tate}}$; these data determine the induced
GRT$_1$ action on the scalar Virasoro line.  Row $\bB$ combines a
weight-filtered associator action with an imaginary-cone comparison
and the Vol.~III compact Hall-source recognition gates.  The value
$K^\kappa=8$ is a target comparison coordinate whose source
interpretation is supplied by those recognition data.
\end{remark}

\section{Higher correction cochains \(\delta^{(n)}\)}
\label{sec:nc-delta-n-higher-explicit}

The symbol \(\delta^{(n)}\) acquires geometric meaning after its source
complex, target complex, differential, and scalar projection have been
specified.  Catalan numbers count planar trees, while motivic
dimensions count classes in a period algebra.  A correction cochain
requires a comparison map from those combinatorial and motivic inputs
to the ordered residue complex.

\begin{proposition}[Construction problem for higher correction cochains]
\label{prop:nc-delta-n-explicit-higher}
\ClaimStatusOpen
\emph{Type signature:}
\textup{(Open quadrant, completed ordered residue presentation,
Beilinson level~\(2\),
\(\mathsf H_{\mathrm{res}}+\mathsf H_{\mathrm{mot}}\)).}
For \(n\ge7\), an explicit correction cochain is determined by the
following data:
\begin{enumerate}[label=\textup{(D\arabic*)}]
\item a represented arity-\(n\) class in the ordered bar complex;
\item the transferred Maurer--Cartan differential and its collision
 strata;
\item inverse Gram forms for every intermediate quasi-primary channel;
\item a scalar contraction compatible with the motivic comparison.
\end{enumerate}
For the Virasoro local model, the first nontrivial inverse Gram factor
is supplied by
\[
 \Lambda={:}TT{:}-\frac{3}{10}\partial^2T,
 \qquad
 \langle\Lambda,\Lambda\rangle=\frac{c(5c+22)}{10}.
\]
These data give an algorithm for \(\delta^{(n)}\).  A closed recurrence
and each polynomial numerator are outputs of that algorithm.
\end{proposition}

\begin{remark}[Planar counts and residue corrections]
\label{rem:nc-delta-n-beyond-leading}
The Catalan number \(C_{n-1}\) records a planar-tree indexing set, and
a motivic dimension \(d_n\) records the size of a source period space.
The ordered residue cochain also contains nonplanar collision channels,
quasi-primary projections, and inverse Gram forms.  The package
\(\mathsf H_{\mathrm{res}}+\mathsf H_{\mathrm{mot}}\) assembles these
ingredients into one typed correction.  Direct computations at
successive arities then determine its poles and numerator polynomials.
\end{remark}
